\documentclass[english,13pt]{article}
%\usepackage{xcolor}
%\usepackage{hyperref}
\usepackage[T1]{fontenc}
\usepackage[latin1]{inputenc}
\usepackage{geometry}
\geometry{verbose,tmargin=3cm,bmargin=3cm,lmargin=2.75cm,rmargin=2.75cm}
\usepackage{float}
\usepackage{mathrsfs}
\usepackage{amsmath}
\usepackage{amssymb}
\usepackage{graphicx}
\usepackage[normalem]{ulem}
%\usepackage{subfigure}
\usepackage{hyperref}
\hypersetup{
    colorlinks=true,
    linkcolor=blue,
    filecolor=magenta,      
    urlcolor=cyan,
    citecolor = red,
}

%\usepackage[notcite,notref]{showkeys}

%\usepackage{subfigure} %this contradicts with hyperref package

%change color of equation number when using \ref
%\hypersetup{
%    colorlinks=true,
%    linkcolor=blue,
%    filecolor=magenta,      
%    urlcolor=cyan,
%    citecolor = blue,
%}

\makeatletter

%%%%%%%%%%%%%%%%%%%%%%%%%%%%%% LyX specific LaTeX commands.
\usepackage{algorithm,algpseudocode}
\newcommand{\algorithmautorefname}{Algorithm}

%\floatstyle{ruled}
%\newfloat{algorithm}{tbp}{loa}
%\providecommand{\algorithmname}{Algorithm}
%\floatname{algorithm}{\protect\algorithmname}

%%%%%%%%%%%%%%%%%%%%%%%%%%%%%% User specified LaTeX commands.
%%%%%%%%%%%%%%%%%%%%%%%%%%%%%% LyX specific LaTeX commands.
%% Because html converters don't know tabularnewline


%%%%%%%%%%%%%%%%%%%%%%%%%%%%%% Textclass specific LaTeX commands


\usepackage{amsthm}


\usepackage{mathrsfs}


\usepackage{amsfonts}


\usepackage{epsfig}


\usepackage{bm}


\usepackage{mathrsfs}


\usepackage{enumerate}


\numberwithin{equation}{section}


\@ifundefined{definecolor}{\@ifundefined{definecolor}
 {\@ifundefined{definecolor}
 {\usepackage{color}}{}
}{}
}{}
%\usepackage{subfigure}

%\newtheorem{theorem}{Theorem}[section]
%\newtheorem{lem}{Lemma}[section]
%\newtheorem{rem}{Remark}[section]
%\newtheorem{prop}{Proposition}[section]
%\newtheorem{cor}{Corollary}[section]
\newtheorem{ass}{Assumption}[section]


\usepackage{subfig}\usepackage[all]{xy}

%\usepackage[auth-sc,affil-it]{authblk}\renewcommand{\baselinestretch}{1.62}
\newcommand{\bbR}{\mathbb R}
\newcommand{\poi}{\widehat{\varphi}}
%\textwidth   5.8in \textheight  9in \oddsidemargin  0.3in
%\evensidemargin 0.3in \topmargin 0in

\newtheorem{theorem}{Theorem}[section]
\newtheorem{lem}{Lemma}[section]
\newtheorem{rem}{Remark}[section]
\newtheorem{prop}{Proposition}[section]
\newtheorem{cor}{Corollary}[section]
\newcounter{hypA}
\newenvironment{hypA}{\refstepcounter{hypA}\begin{itemize}
  \item[({\bf A\arabic{hypA}})]}{\end{itemize}}
\newcounter{hypB}
\newenvironment{hypB}{\refstepcounter{hypB}\begin{itemize}
 \item[({\bf H\arabic{hypB}})]}{\end{itemize}}
\newcounter{hypD}
\newenvironment{hypD}{\refstepcounter{hypD}\begin{itemize}
 \item[({\bf D\arabic{hypD}})]}{\end{itemize}}
\newcounter{hypW}
\newenvironment{hypW}{\refstepcounter{hypW}\begin{itemize}
  \item[({\bf W\arabic{hypW}})]}{\end{itemize}}


\newcommand{\Exp}{\mathbb{E}}
\newcommand{\Con}{X_{l_d(t_{k-1}(d)),j}^{\prime,i}}


\usepackage{babel}\date{}



\usepackage{babel}

\makeatother

\usepackage{babel}

\newcommand{\calA}{\mathcal{A}}
\newcommand{\cO}{\mathcal{O}}
\newcommand{\calB}{\mathcal{B}}
\newcommand{\calC}{\mathcal{C}}
\newcommand{\calD}{\mathcal{D}}
\newcommand{\calE}{\mathcal{E}}
\newcommand{\calF}{\mathcal{F}}
\newcommand{\calG}{\mathcal{G}}
\newcommand{\calH}{\mathcal{H}}
\newcommand{\calI}{\mathcal{I}}
\newcommand{\calJ}{\mathcal{J}}
\newcommand{\calK}{\mathcal{K}}
\newcommand{\calL}{\mathcal{L}}
\newcommand{\calM}{\mathcal{M}}
\newcommand{\calN}{\mathcal{N}}
\newcommand{\calO}{\mathcal{O}}
\newcommand{\calP}{\mathcal{P}}
\newcommand{\calQ}{\mathcal{Q}}
\newcommand{\calR}{\mathcal{R}}
\newcommand{\calS}{\mathcal{S}}
\newcommand{\calT}{\mathcal{T}}
\newcommand{\calU}{\mathcal{U}}
\newcommand{\calV}{\mathcal{V}}
\newcommand{\calW}{\mathcal{W}}
\newcommand{\calX}{\mathcal{X}}
\newcommand{\calY}{\mathcal{Y}}
\newcommand{\calZ}{\mathcal{Z}}
\newcommand{\cD}{\mathcal{D}}
\newcommand{\cA}{\mathcal{A}}
\newcommand{\cK}{\mathcal{K}}
\newcommand{\cQ}{\mathcal{Q}}
\newcommand{\cN}{\mathcal{N}}
\newcommand{\cF}{\mathcal{F}}
\newcommand{\cE}{\mathcal{E}}
\newcommand{\cV}{\mathcal{V}}
\newcommand{\cM}{\mathcal{M}}
\newcommand{\cP}{\mathcal{P}}
\newcommand{\sX}{\mathsf X}
\newcommand{\sY}{\mathsf Y}


\newcommand{\bbA}{\mathbb{A}}
\newcommand{\bbB}{\mathbb{B}}
\newcommand{\bbC}{\mathbb{C}}
\newcommand{\bbD}{\mathbb{D}}
\newcommand{\bbE}{\mathbb{E}}
\newcommand{\bbF}{\mathbb{F}}
\newcommand{\bbG}{\mathbb{G}}
\newcommand{\bbH}{\mathbb{H}}
\newcommand{\bbI}{\mathbb{I}}
\newcommand{\bbJ}{\mathbb{J}}
\newcommand{\bbK}{\mathbb{K}}
\newcommand{\bbL}{\mathbb{L}}
\newcommand{\bbM}{\mathbb{M}}
\newcommand{\bbN}{\mathbb{N}}
\newcommand{\bbO}{\mathbb{O}}
\newcommand{\bbP}{\mathbb{P}}
\newcommand{\bbQ}{\mathbb{Q}}
%\newcommand{\bbR}{\mathbb{R}}
\newcommand{\bbS}{\mathbb{S}}
\newcommand{\bbT}{\mathbb{T}}
\newcommand{\bbU}{\mathbb{U}}
\newcommand{\bbV}{\mathbb{V}}
\newcommand{\bbW}{\mathbb{W}}
\newcommand{\bbX}{\mathbb{X}}
\newcommand{\bbY}{\mathbb{Y}}
\newcommand{\bbZ}{\mathbb{Z}}

\def\QQuad{\hskip10\fontdimen16\font}

\newcommand{\KL}[1]{{\color{blue}[KL: \textbf{#1}]}} % Sara  

\newcommand{\NK}[1]{{\color{blue} {#1}}} 

\begin{document}

%+Title
\begin{center}

{\Large \textbf{Particle Filtering for a Class of State-Space Models with Low and Degenerate Observational Noise}}

\vspace{0.5cm}

BY  ABYLAY ZHUMEKENOV$^{1}$,  ALEXANDROS BESKOS$^{2}$,  DAN CRISAN$^{3}$, AJAY JASRA$^{1}$ \& NIKOLAS KANTAS$^{3}$

{\footnotesize $^{1}$School of Data Science,  The Chinese University of Hong Kong,  Shenzhen,  Shenzhen, CN.}\\
{\footnotesize $^{2}$Department of Statistical Science, University College London, London, WC1E 6BT, UK}\\
{\footnotesize $^{3}$Department of Mathematics, Imperial College London, London, SW7 2AZ, UK}\\
{\footnotesize E-Mail:\,} \texttt{\emph{\footnotesize  abylayzhumekenov@cuhk.edu.cn; 
a.beskos@ucl.ac.uk; d.crisan@ic.ac.uk;  
ajayjasra@cuhk.edu.cn; n.kantas@ic.ac.uk
}}

\end{center}



\begin{abstract}
We consider the discrete-time filtering problem in scenarios where the observation noise is  low or degenerate.
We focus on the case where the observation equation is a linear function of the state and the data   involve additive noise.  However,  we place minimal assumptions on the hidden state process.  For such a class of models  we derive new particle filters (PFs) with the key property that their performance is robust to the size of the observation noise. As a consequence,  the developed PFs are well-defined in the limiting case of degenerate observation noise.  Indicatively,  we prove (under assumptions) that the PF applied  in this low noise setting inherits the properties of the PF used in the degenerate case.
  We extend our framework to the case where the hidden states are drawn from a diffusion process.  In this scenario we develop new PFs which are robust to both low noise and fine levels of time discretization.  We illustrate our algorithms numerically on several examples.
\\
\textbf{Keywords}:  Particle filters,  low \& degenerate noise,  diffusions,  time discretization.
\end{abstract}

\section{Introduction}

The discrete-time filtering problem considers an unobserved Markov chain $\{X_n\}_{n\geq 0}$ (the signal or hidden chain) and associated data
$\{Y_n\}_{n\geq 1}$.  The data at any time $k$ are assumed to depend on   $\{X_n\}_{n\geq 0}$ only via the position $X_k$ of the hidden state.  The objective of filtering is to compute expectations with respect to (w.r.t.) the sequence of \emph{filtering distributions}, i.e.~the conditional distribution of the signal at each time $n$,  given all the data observed up-to that time, sequentially in the time index. 
The filtering problem has numerous applications in statistics, applied mathematics, finance and beyond;  we refer the reader to the texts \cite{bain,cappe,delmoral} and the references therein for a more complete summary.

In most cases of practical interest  the filter is unavailable analytically and one must resort to numerical approximations.  A plethora of relevant algorithms have been developed,  see e.g.~\cite{bain,cappe,delmoral} for some coverage of the subject.  In this article we focus on the approach called particle filters (PFs).  These are Monte Carlo,  simulation-based methods that simulate a collection of $N\in\mathbb{N}$ samples in parallel and undergo importance sampling and resampling operations. Such methods are very well understood from a mathematical perspective (e.g.~\cite{delmoral}) and provide almost surely consistent approximations of the filter as $N\rightarrow\infty$.  PFs are typically effective when the dimension of the hidden state, $d_x$,  is moderate,  although there are several recent attempts to further extend the methodology to cover some special-case high-dimensional problems (e.g.~\cite{pf_ns,pf_sns}).

In this article we focus on the case when the observations have low or degenerate noise.  For instance, one could consider the case:
%
\begin{equation}
\label{eq:add_noise_ex}
Y_n = A_n X_n + \Delta^{1/2}\epsilon_n,
\end{equation}
%
where $(Y_n,X_n)\in\mathbb{R}^{d_y}\times\mathbb{R}^{d_x}$,  $\Delta\geq 0$, 
$\{\epsilon_n\}_{n\geq 1}$ is a sequence of $\mathbb{R}^{d_y}$-valued independent and identically distributed (i.i.d.) noises and $A_n$ a sequence of $d_y\times d_x$ real matrices.
 of full row rank.
For technical reasons that will become apparent in the sequel, we assume that $1\leq d_y<d_x$ and that $A_n$ is of full row rank for every $n\ge 1$.
The low noise case could be interpreted as $\Delta$ small (e.g.~$\Delta=10^{-6}$) whereas the  degenerate case corresponds to $\Delta=0$.  As noticed by \cite{zhu},  such a setting is
important in many applications  where it is expected that the signal dynamics is strongly driving the observed data; see e.g.~\cite{app1,app2}.  From a mathematical perspective,  in the case $\Delta=0$ one has to define the filter on the following sequence of \emph{manifolds}:
%
\begin{equation*}
\mathsf{M}^\star_n=\{x\in\mathbb{R}^{d_x}:y_n-A_nx=0\}.
\end{equation*}
%
 Indeed this setting has been explored,  for instance,  in \cite{zhu}.  In a non-sequential context when one targets a single distribution (e.g., the smoothing law of $X_0,\dots,X_n|y_1,\dots,y_n$ with $n$ fixed), sampling on manifolds has been quite well investigated using Markov chain Monte Carlo (MCMC) in several interesting contexts; see for example
\cite{bhar,byrne,graham_st,graham,zappa}. The extension to the filtering problem seems only to be considered in \cite{zhang,zhu},  although \cite{zhang} do not consider a manifold that has been defined by the data such as $\mathsf{M}_n$ written above.  The approach of \cite{zhu} is to leverage sequential MCMC and as such it relies on the afore-mentioned MCMC literature.  Such methodology is particularly relevant in high-dimensions ($d_x$ large) but can be computationally expensive for moderate dimensional problems. Instead, within this latter class of filtering problems it is natural to consider PFs and this is the direction we follow in the present work.

 In a more general context, understanding the filtering problem in this low noise regime has been the focus of many studies investigating continuous time dynamics and devising robust (extended) Kalman-Bucy type approximations of the filter \cite{brigo,katzur, milheiro, picard1, picard2} or establishing large deviations results for low $\Delta$  \cite{hijab,pardoux}, see also \cite{kutoyants} for a recent review. In general, we need to stress that the filtering distributions are typically analytically intractable except in fairly simple models,  such as ones involving linear and Gaussian dynamics.  Using approximations such as the extended Kalman filter or variants are biased and most often do not perform well for many models. Given also the increasing accuracy of many sensors used in practice it is important to design PFs that perform well for degenerate or low observation noise case.

The application of {plain} PFs for low noise problems seldom works well due to the informative likelihoods resulting to very low weights for most particles. This is well-known and illustrated later in this article. The main idea of our work is to {reparameterize} and cast the given filtering problem as one of sampling on a sequence of manifolds, \emph{both in the degenerate setting and in the low noise one}.  
{For instance},  in the context of \eqref{eq:add_noise_ex}, for $\Delta > 0$, we work with the sequence of manifolds: 
%
\begin{equation*}
\widetilde{\mathsf{M}}_n=\{(x,\epsilon)\in\mathbb{R}^{d_x}\times\mathbb{R}^{d_y}:y_n-A_nx-\Delta^{1/2}\epsilon=0\}.
\end{equation*} 
%
A key idea then is that one can carefully define PFs on $\mathsf{M}_n$ and $\widetilde{\mathsf{M}}_n$, so that the algorithm in the low noise case will be robust by admitting a cost $\mathcal{O}(1)$  w.r.t.~$\Delta$, and will thus converge, in a proper sense, to the PF for the degenerate case of $\Delta=0$. 

 In this paper we explore methodologies with the above properties in the presence of linear observation models such as 
\eqref{eq:add_noise_ex}. Our contributions can be summarized as follows:
%
\begin{itemize}
\item[(i)]We derive new PFs for cases of low and degenerate noise, within the setting of linear observation models.  We place minimal assumptions on the hidden process $\{X_n\}_{n\geq 0}$.
\item[(ii)]In the scenario of (i) and under assumptions, we prove that the $\mathbb{L}_r$-error (for a positive integer $r\ge 1$) of the PF in the low noise case will converge, as $\Delta\downarrow 0$, to the $\mathbb{L}_r$-error of the PF in the degenerate noise case.
\item[(iii)] For additive Gaussian errors both in the hidden chain and the observations, we derive the optimal proposal and we investigate 
the corresponding PFs both under low and degenerate noise case. In particular, we show that the weights of the PF for $\Delta>0$ converge  to the ones for $\Delta=0$ as $\Delta\downarrow 0$.
\item[(iv)]In the case where the hidden dynamics follow a {continuous time} diffusion process
which is observed at discrete times,  
we derive a new PF that is robust not only to low noise but also to fine time discretization error. This appears to be the first PF to have this property.
\end{itemize}
%
The significance of this work is not only restricted to filtering.  The new PFs in this article can be used for parameter estimation
associated to diffusion models,  such as was done in \cite{diff_bridge,non_synch} using the multilevel Monte Carlo method.  We note in terms of the fourth contribution,  that \cite{bierkens} develop a methodology for the case of hypoelliptic diffusions and degenerate noise (with linear constraints),  stating that PFs can be used in their context.  However, they do not make the point that this enables the low noise case,  which is one of the main findings of this work.

This article is structured as follows.  In Section \ref{sec:model} we present the collection of models we will consider.
In Section \ref{sec:case_study} we give two particular types of models with PF algorithms and mathematical results.
In Section \ref{sec:diff} we show how to extend our methodology to the case where the hidden dynamics follow a diffusion process.
In Section \ref{sec:numerics} we present our numerical results.  The proofs of our mathematical results can be found in Appendix \ref{app:proof}.

\section{Modeling and Filtering}\label{sec:model}

\subsection{State-Space Model}

In this section, we work in a general setting in terms of the class of models we consider. 
Our objective is to present what we call the 
`noise parameterization' (as an alternative to the `natural parametrization') and the degenerate noise case in full generality.  As we stated in the introduction,  later on we are only able to develop methodology for linear observation regimes; this is explained in the next section.  
%We note,  however,  we will present the models in a more complete framework than the linear observation case. 

We consider the following state-space model (SSM) for $n\in\mathbb{N}$ and $X_0\in\mathbb{R}^{d_x}$:
%
\begin{align*}
Y_n & = H_n(X_n,\epsilon_n), \\
X_n & =  F_n(X_{n-1},\nu_n),
\end{align*}
%
where we have that: $(Y_n,X_n)\in\mathbb{R}^{d_y+d_x}$, $(d_x,d_y)\in\mathbb{N}^2$;  $\epsilon_n\in\mathbb{R}^{d_x}$, $\nu_n\in\mathbb{R}^{d_y}$ form  sequences of noises that are  independent both across $n$ and of each other, with all variables having  positive and smooth Lebesgue densities;  $H_n:\mathbb{R}^{d_x+d_y}\rightarrow\mathbb{R}^{d_y}$, $F_n:\mathbb{R}^{2d_x}\rightarrow\mathbb{R}^{d_x}$ are  sequences of  continuously differentiable functions with Jacobians that are of full row rank.

\subsection{Natural Parametrization}

In the given scenario we can define the filtering densities w.r.t.~the $d_x$-dimensional Lebesgue measure as follows.
Denoting the positive conditional Lebesgue densities of $Y_n|X_n$ and $X_n|X_{n-1}$ as $h_n(y_n|x_n)$ and $f_n(x_n|x_{n-1})$ respectively, we have that the filter (see e.g.~\cite{bain} and the references therein) at time $1$ writes as follows (with some abuse of notation, here and elsewhere below we use the same symbol both for a distribution and for its density w.r.t.~a reference measure implied by the context): 
%
\begin{align}
\label{eq:f-1}
\pi_1(x_1)\equiv \frac{d\pi_1}{d\lambda_{d_x}}(x_1) =  \frac{h_1(y_1|x_1)f_1(x_1|x_0)}{\int_{\mathbb{R}^{d_x}}h_1(y_1|x_1)f_1(x_1|x_0)\lambda_{d_x}(dx_1)},
\end{align}
%
where $\lambda_{d_x}$ is the $d_x$-dimensional Lebesgue measure and we supress the dependence of $y_1$ in the expression for the filtering density $\pi_1(x_1)$.
The filter at any subsequent time point $n\geq 1$ has the following prediction-updating formula:
%
\begin{align}
\label{eq:f-n}
\pi_n(x_n) \equiv   \frac{d\pi_n}{d\lambda_{d_x}}(x_n) =  \frac{h_n(y_n|x_n)\int_{\mathbb{R}^{d_x}} f_n(x_n|x_{n-1})\pi_{n-1}(x_{n-1})\lambda_{d_x}(dx_{n-1})}{\int_{\mathbb{R}^{d_x}}h_n(y_n|x_n)\int_{\mathbb{R}^{d_x}} f_n(x_n|x_{n-1})\pi_{n-1}(x_{n-1})\lambda_{d_x}(dx_{n-1})\lambda_{d_x}(dx_n)}.
\end{align}
%

In practice one cannot compute the filter  and one must resort to numerical methods. In standard settings, PFs are well known to provide accurate approximations of $\pi_n(x_n)$ and this is for an arbitrary long $n$ under favourable stability assumptions for $F_n$ and non-degeneracy of $\epsilon_n$. A PF algorithm is in fact designed to target at time $n$ the smoothing density
$$
\pi_n({x}_{1:n}) \propto \prod_{k=1}^n h_k(y_k|x_k)
f_k({x}_k|x_{k-1}),
$$
from which $\pi_n(x_n)$ can be derived as its final marginal. It should be clear that in this formulation one cannot deal directly with a degenerate case for the conditional likelihood $h_k$ and we will redefine an equivalent problem on an appropriate space. In addition, for the low noise case one should note that the region where $h_k$ has substantial value will be quite small and becomes a low probability event for traditional sampling proposals to reach it and hence standard PFs can be ineffective.
%Consider when $d_y=1$ and $\epsilon_n\sim\mathcal{N}(0,\sigma_n^2)$
%with $\mathcal{N}(0,\sigma_n^2)$ the one-dimensional normal distribution of mean zero and variance $\sigma_n^2$ very small
%say $10^{-6}$ we call this a special case of low observational noise.  
We proceed below with a different parametrization for the same SSM model which will allow for the development of effective numerical methods in the scenario of low noise for the observations.  %it can be difficult to design an efficient 
%numerical methods and to help deal with such a scenario we will consider a SSM under a noise-based parameterization which we now detail.

% a particle filter (PF).  A main ingredient of this methodology is the (positive) proposal density $q_n(x_n|x_{n-1})$ which defines an un-normalized incremental weight (UIW):
%$$
%w_n(x_{n-1:n}) = \frac{h_n(y_n|x_n)f_n(x_n|x_{n-1})}{q_n(x_n|x_{n-1})}.
%$$
%Two examples are the bootstrap particle filter with proposal $f_n(x_n|x_{n-1})$ and UIW
%$$
%w_n(x_n) = h_n(y_n|x_n)
%$$
%and the optimal proposal
%$$
%q_n(x_n|x_{n-1}) = \frac{ h_n(y_n|x_n)f_n(x_n|x_{n-1}) }{\int_{\mathbb{R}^{d_x}}
%h_n(y_n|x_n)f_n(x_n|x_{n-1})\lambda(dx_n) }
%$$
%with UIW
%$$
%w_{n}(x_{n-1}) = \int_{\mathbb{R}^{d_x}}h_n(y_n|x_n)f_n(x_n|x_{n-1})\lambda(dx_n).
%$$
%In many cases the optimal proposal is not available, but there are special cases where it can be sampled and computed.  We call a PF which works in the scenario just described as a PF associated to a SSM under the natural parameterization. 
%Consider when $d_y=1$ and $\epsilon_n\sim\mathcal{N}(0,\sigma_n^2)$
%with $\mathcal{N}(0,\sigma_n^2)$ the one-dimensional normal distribution of mean zero and variance $\sigma_n^2$ very small
%say $10^{-6}$ we call this a special case of low observational noise.  In such cases it can be difficult to design an efficient particle filter as UIW could explode and  (simultanesouly) have exploding variance.
%To deal with such as scenario we will consider a SSM under a noise-based parameterization which we now detail.


\subsubsection{Particle Filter for the Natural Parameterization}

We will present a generic PF for this parameterization when $\Delta>0$.  We will use the ancestor representation of the resampling indices denoted as $({a}_n^1,\dots,{a}_n^N)\in\{1,\dots,N\}^N$ and we use the convention that ${a}_0^i=i$ for each $i\in\{1,\dots,N\}$.  %For the initialization  we will also write ${x}_{0}^{{a}_{0}^j}=x_0$, $a_0^i=i$ for $i\in\{1,\dots,n\}$ and. 
For the importance sampling part, $q_n$ is a sequence of positive Markov transition densities on $\mathbb{R}^{d_x}$ and for the weights we set for any $(x',x)\in\mathbb{R}^{2d_x}$. 
$$
w_n(x',x) = \frac{h_n(y_n|x)f_n(x|x')}{q_n(x|x')}.
$$
Note that ${q}_n$ can depend on $\Delta$ and when $n=1$ we may (or may not) condition on $x_0$
in ${q}_1$ but write ${q}_1(\cdot|x_0)$ and use $x_0^i=x_0$ for all $i$.
%The PF is then the following.

%For this problem we already know the PF for the filter in the noise parameterization case as well as the degenerate noise case.
% For completeness,  we detail the natural parameterization case when $\Delta>0$. 
Let $\text{Cat}(\cdot;k,\mathbf{p})$ denote the categorical distribution with $k>0$ categories and $\mathbf{p}=(p_1,\ldots,p_k)$ consist a vector of event probabilies such that $p_i\geq 0$ and $\sum_{j=1}^k p_j =1$. Then the PF algorithm proceeds as follows:
\begin{itemize}
\item{Initialize: For $i\in\{1,\dots,N\}$,  sample $X_1^i|x_0$ from $q_1(\cdot|x_0)$. Set $n=2$.}
\item{Iterate: For $i\in\{1,\dots,N\}$,  sample $a_{n-1}^i\sim\text{Cat}(\cdot;N,\mathbf{w}_{n-1})$ with
$$
\mathbf{w}_{n-1}=\left( \frac{w_{n-1}(x_{n-2}^{a_{n-2}^1},x_{n-1}^{1})}
{\sum_{s=1}^N w_{n-1}(x_{n-2}^{a_{n-2}^s},x_{n-1}^s)},\dots,\frac{w_{n-1}(x_{n-2}^{a_{n-2}^N},x_{n-1}^{N})}
{\sum_{s=1}^N w_{n-1}(x_{n-2}^{a_{n-2}^s},x_{n-1}^s)}\right).
%\frac{w_{n-1}(x_{n-2}^{a_{n-2}^i},x_{n-1}^{a_{n-1}^i})}
%{\sum_{s=1}^N w_{n-1}(x_{n-2}^{a_{n-2}^s},x_{n-1}^s)}.
$$
% $a_{n-1}^i|x_{n-2}^{1:N},x_{n-1}^{1:N}$ i.i.d.~from
For $i\in\{1,\dots,N\}$,  sample $X_n^i|x_{n-1}^{a_{n-1}^i}$ independently from $q_n(\cdot|x_{n-1}^{a_{n-1}^i})$.
Set $n=n+1$ and return to the start of the bullet point.}
\end{itemize}
The optimal proposal in terms of minimizing the variance of the weights conditional on ${x}_{n-1}$ is 
$$
q_n^{\text{opt}}(x_n|x_{n-1}) = \frac{h_n(y_n|x_n)f_n(x_n|x_{n-1})}
{\int_{\mathbb{R}^{d_x}} h_n(y_n|x_n)f_n(x_n|x_{n-1})\lambda_{d_x}(dx_n)}.
$$
Often used (and maligned) is the \emph{bootstrap PF} where one uses $q_n(x|x')=f_n(x|x')$.


\subsection{Noise Parametrization and Degenerate Noise}

\subsubsection{Noise Parametrization}

We set $c_n(x_n,\epsilon_n) = y_n-H_n(x_n,\epsilon_n)$ and define the manifold:
% 
\begin{equation}
\label{eq:M-general-low}
 \widetilde{\mathsf{M}}_n=\{(x_n,\epsilon_n)\in\mathbb{R}^{d_x}\times\mathbb{R}^{d_y}:
c_n(x_n,\epsilon_n) = 0\} .
\end{equation}
%
Conditionally on $x_{n-1}$ (but not on $y_n$), the variables $(x_n,\epsilon_n)$ have a Lebesgue density: 
%
\begin{equation*}
p_n(x_{n}, \epsilon_n)  = f_n(x_n|x_{n-1})\,p_n(\epsilon_n) .
\end{equation*}
% 
Here $p_n(\epsilon_n)$ is the Lebesgue density of $\epsilon_n$.
It is useful to think  of $p_n(x_n,\epsilon_n)$ as the prior density for $(x_n, \epsilon_n)$ given $x_{n-1}$ before consideration of $y_n$.  Once $y_n$ is received, $(x_n,\epsilon_n)$ will lie on the manifold  $\widetilde{\mathsf{M}}_n$ in (\ref{eq:M-general-low}). In particular, the posterior distribution of $(x_n,\epsilon_n)$ is their prior with its support now restricted on  $\widetilde{\mathsf{M}}_n$ (up-to a constant and Jacobian factor). One can obtain an expression for the density of the posterior upon consideration of the appropriate reference measure on the manifold.  

In detail, we assume that the ambient space $\mathbb{R}^{d_x+d_y}$ is equipped with a metric tensor with a fixed positive definite matrix representation $M$. Then, it is simple to show that the filter density 
of $(x_n,\epsilon_n)$ w.r.t.~the Riemannian measure on $\widetilde{\mathsf{M}}_n$ under $M$, denoted $\mu_{\widetilde{\mathsf{M}}_n}$,
is as follows (see e.g.~\cite{zhu}).  We set:
%
\begin{align}
\label{eq:Gram-e}
\widetilde{G}_n(x_n,\epsilon_n) = \textrm{det}(\partial c_n(x_n,\epsilon_n)M^{-1}\partial c_n(x_n,\epsilon_n)^{\top})^{-1/2},
\end{align}
%
where $\partial c_n$ is the Jacobian matrix and $\textrm{det}(\cdot)$ is the determinant.
For time 1 we have the filtering density:
%
\begin{align*}
\widetilde{\pi}_1(x_1,\epsilon_1) \equiv \frac{d\widetilde{\pi}_1}{d\mu_{\widetilde{\mathsf{M}}_1}}(x_1,\epsilon_1)  =   \frac{
\widetilde{G}_1(x_1,\epsilon_1)f_1(x_1|x_0)p_1(\epsilon_1)}{
\int_{\widetilde{\mathsf{M}}_1}
\widetilde{G}_1(x_1,\epsilon_1)f_1(x_1|x_0)p_1(\epsilon_1)
\mu_{\widetilde{\mathsf{M}}_1}(d(x_n,\epsilon_n))
}.
\end{align*}
%
At any subsequent time $n\geq 2$, we have:
%
\begin{align*}
&\widetilde{\pi}_n(x_n,\epsilon_n) \equiv \frac{d\widetilde{\pi}_n}{d\mu_{\widetilde{\mathsf{M}}_n}}(x_n,\epsilon_n)    \\ 
 &\qquad = \frac{
\widetilde{G}_n(x_n,\epsilon_n)p_n(\epsilon_n) \int_{\widetilde{\mathsf{M}}_{n-1} } f_n(x_n|x_{n-1})
\widetilde{\pi}_n(x_{n-1},\epsilon_{n-1}) \mu_{\widetilde{\mathsf{M}}_{n-1}}(d(x_{n-1},\epsilon_{n-1}))
}{
\int_{\widetilde{\mathsf{M}}_n}
\widetilde{G}_n(x_n,\epsilon_n)p_n(\epsilon_n) \int_{\widetilde{\mathsf{M}}_{n-1} } f_n(x_n|x_{n-1})
\widetilde{\pi}_n(x_{n-1},\epsilon_{n-1}) \mu_{\widetilde{\mathsf{M}}_{n-1}}(d(x_{n-1},\epsilon_{n-1}))
\mu_{\widetilde{\mathsf{M}}_n}(d(x_n,\epsilon_n))
}.
\end{align*}
%
We note that under the law $d\widetilde{\pi}_n(x_n, \epsilon_n)$, $n\in \mathbb{N}$,  the marginal  of $x_n$ is as given in  expressions (\ref{eq:f-1})--(\ref{eq:f-n}) for the natural parametrization. As before one can also write a path smoothing density
\begin{align}
\widetilde{\pi}_n(x_{0:n},\epsilon_{0:n}) \equiv \frac{d\widetilde{\pi}_n}{d\mu_{\bigotimes_{k=1}^n \widetilde{\mathsf{M}}_k}}(x_{0:n},\epsilon_{0:n})    \propto 
\prod_{k=1}^n \widetilde{G}_k(x_k,\epsilon_k)p_k(\epsilon_k)  f_k(x_k|x_{k-1}),
\label{eq:smoothing_M_n_low_noise}
\end{align}
which will form the target density of a PF.

We remark that one can obtain some corresponding re-parametrized expressions by working  with the noise variables $(\nu_n)$ of the hidden state instead of $(x_n)$. However, this is not useful in our setting  as we focus on observational noise.  As we will establish in the next section, when the noise in the observation equation is low,  one can consider numerical algorithms under the noise parametrization and ultimately,  as the size of the noise goes to $0$,  recover an algorithm associated to the degenerate noise filtering problem.  For this latter problem, we obtain in the subsection that follows expressions analogous to the ones in the current subsection. 


%However,  the issues would be that (i) would have to consider a PF on path space,  which seldom works well and (ii) even if the PF works well it would have a cost that is quadratic in time and (iii) if the noises in the past are also updated one has to verify the constraints are also satisfied.

%In this parameterization one can consider a PF with proposal transition kernel:
%$$
%Q_n(d(x_n,\epsilon_n)|x_{n-1},\epsilon_{n-1}) = q_n(x_n,\epsilon_n|x_{n-1},\epsilon_{n-1})
%\mu^M_{\mathsf{M}_n}(d(x_n,\epsilon_n)).
%$$
%For instance the optimal proposal would be
%$$
%Q_n(d(x_n,\epsilon_n)|x_{n-1},\epsilon_{n-1}) =
%\frac{
%G_n(x_n,\epsilon_n)p_n(\epsilon_n)f(x_n|x_{n-1})\mu^M_{\mathsf{M}_n}(d(x_n,\epsilon_n))
%}
%{
%\int_{\mathsf{M}_n} G_n(x_n,\epsilon_n)p_n(\epsilon_n) f(x_n|x_{n-1})\mu^M_{\mathsf{M}_n}(d(x_n,\epsilon_n))
%}
%$$
%yielding an UIW
%$$
%w_n(x_{n-1}) = \int_{\mathsf{M}_n} G_n(x_n,\epsilon_n)p_n(\epsilon_n) f(x_n|x_{n-1})\mu^M_{\mathsf{M}_n}(d(x_n,\epsilon_n)).
%$$
%In many problems,  such a proposal will not be possible to be used;  we will consider a special case in the next section.
%A general design of proposals is considered later on.  

\subsubsection{Degenerate Noise}

We now consider degenerate observation noise and show the steps for obtaining expression for the sequence of filtering densities w.r.t.~appropriate reference measures.  We write the SSM under consideration as:
%
\begin{align*}
Y_n & = H_n^{\star}(X_n), \\
X_n & =  F_n(X_{n-1},\nu_n),
\end{align*}
%
where: $(Y_n,X_n)\in\mathbb{R}^{d_y+d_x}$, $(d_x,d_y)\in\mathbb{N}^2$, $d_y<d_x$;  $\nu_n\in\mathbb{R}^{d_x}$ is a sequence of independent noises (across $n$) with positive and smooth Lebesgue density; and $H_n^{\star}:\mathbb{R}^{d_x}\rightarrow\mathbb{R}^{d_y}$, $F_n:\mathbb{R}^{2d_x}\rightarrow\mathbb{R}^{d_x}$ are sequences of  continuously differentiable functions whose Jacobian is of full row rank.  Note that we require $d_y<d_x$ as we aim to work with an embedded manifold.

As above, we set $c_n(x_n) = y_n-H_n^{\star}(x_n)$ and define the manifold:
%
\begin{align}
\label{eq:M-general-zero}
 \mathsf{M}_n^{\star}=\{x_n\in\mathbb{R}^{d_x}:
c_n(x_n) = 0\}.
\end{align}
%
We assume that $\mathbb{R}^{d_x}$ is equipped with a metric tensor with a fixed positive definite matrix representation, denoted $M$.  We set:
%
\begin{equation}
\label{eq:Gram}
G_n(x_n) = \textrm{det}(\partial c_n(x_n)M^{-1}\partial c_n(x_n)^{\top})^{-1/2}.
\end{equation}
%
As before the filter densities w.r.t.~$\mu_{\mathsf{M}_n^{\star}}$ can be written down as
%
\begin{align*}
\pi_1^{\star}(x_1) \equiv 
\frac{d\pi_1^{\star}}{d\mu_{\mathsf{M}_{1}^{\star}}}(x_1) = \frac{
G_1(x_1)f_1(x_1|x_0)}{
\int_{\mathsf{M}_1^{\star}}
G_1(x_1)f_1(x_1|x_0)
\mu_{\mathsf{M}_1^{\star}}(dx_1)
}.
\end{align*}
%
At any subsequent time $n\geq 2$, we have: 
%
\begin{align*}
\pi_n^{\star}(x_n) = \frac{d\pi_n^{\star}}{d\mu_{\mathsf{M}_{n}^{\star}}}(x_n) =   \frac{
G_n(x_n) \int_{\mathsf{M}_{n-1}^{\star} } f_{n}(x_n|x_{n-1})
\pi_{n-1}^{\star}(x_{n-1}) \mu_{\mathsf{M}_{n-1}^{\star}}(dx_{n-1})
}{
\int_{\mathsf{M}_n^{\star}}
G_n(x_n)\int_{\mathsf{M}_{n-1}^{\star} } f_n(x_n|x_{n-1})
\pi_{n-1}^{\star}(x_{n-1}) \mu_{\mathsf{M}_{n-1}^{\star}}(dx_{n-1})
\mu_{\mathsf{M}_n^{\star}}(dx_n)
}.
\end{align*}
%For a PF,  the optimal proposal would be
%$$
%Q_n(dx_n|x_{n-1}) =
%\frac{
%\tilde{G}_n(x_n)f(x_n|x_{n-1})\mu^M_{\mathsf{\tilde{M}}_n}(dx_n)
%}
%{
%\int_{\mathsf{\tilde{M}}_n} \tilde{G}_n(x_n) f(x_n|x_{n-1})\mu^M_{\mathsf{\tilde{M}}_n}(dx_n)
%}
%$$
%yielding an UIW
%$$
%w_n(x_{n-1}) = 
%\int_{\mathsf{\tilde{M}}_n} \tilde{G}_n(x_n) f(x_n|x_{n-1})\mu^M_{\mathsf{\tilde{M}}_n}(dx_n).
%$$
Similarly for smoothing density we have
\begin{align}
{\pi}_n^{\star}(x_{0:n}) \equiv \frac{d\widetilde{\pi}_n}{d\mu_{\bigotimes_{k=1}^n \mathsf{M}^{\star}_k}}(x_{0:n})    \propto 
\prod_{k=1}^n {G}_k(x_k) f_k(x_k|x_{k-1}).
\label{eq:smoothing_M_n_degenerate}
\end{align}


\section{Filtering with Linear Observation and  Additive Noise}\label{sec:case_study}

In this section we consider the case of general additive noise (Section \ref{sec:noise_add}) and derive PFs in the low and degenerate noise cases.  We then examine (Section \ref{sec:spec_add}) the case where the hidden process also involves additive Gaussian noise, similarly to the observation process, and derive in such a setting the optimal proposals for the low and degenerate noise models.

\subsection{Reparameterizations and Particle Filters}\label{sec:noise_add}

We consider the SSM given as follows,  for $n\in\mathbb{N}$ and $X_0\in\mathbb{R}^{d_x}$:
%
\begin{align}
\label{eq:obs-linear}
Y_n  = A_n X_n + \Delta^{1/2}\epsilon_n ,
\end{align}
%
where  $Y_n\in\mathbb{R}^{d_y}$,  $1\leq d_y < d_x$,  $A_n\in\mathbb{R}^{d_y\times d_x}$ is of full row rank,  $\epsilon_n\in\mathbb{R}^{d_y}$ forms a sequence of independent noises with positive Lebesgue density $p_n(\cdot)$
and $\Delta\geq 0$.  For the hidden dynamics we do not impose any strong constraints beyond the existence of a  positive Lebesgue transition density which can be evaluated pointwise; it is written as $f_n(x_n|x_{n-1})$ as before. 
\begin{rem}
In this article we present noise degeneracy as the case that in the observation equation,  as $\Delta\downarrow 0$ there is no noise at all.  However,  one could consider,  say $Y_n\in\mathbb{R}^3$,  and an observation equation of the type
$$
Y_n = A_n X_n + \mathtt{P}(\Delta)\epsilon_n,
$$
where we define
$$
\mathtt{P}(\Delta)=
\begin{bmatrix}
1 & 0 & 0\\
1 & 1 & 0 \\
\Delta & \Delta & \Delta
\end{bmatrix}.  
$$
In this case one could send $\Delta\downarrow 0$ which corresponds to removing the observation noise in the third observation component.  This type of scenario could be dealt with in the forthcoming commentary.
This is
in the sense that one would define a filter (noise parameterization) and a corresponding algorithm in the low noise case and the case $\Delta=0$ and that there is some convergence of the low-noise case to this limiting and degenerate case.  The main changes would be mainly notational and hence, for ease of exposition,  we keep the simple scenario that we have started with.
\end{rem}

Recall our notation for the manifolds: 
%
\begin{align*}
 \widetilde{\mathsf{M}}_n&=\{(x_n,\epsilon_n)\in\mathbb{R}^{d_x+d_y}:
y_n - A_n x_n - \Delta^{1/2}\epsilon_n = 0\}; \\ 
 \mathsf{M}_n^{\star}  &=\{x_n\in\mathbb{R}^{d_x}:
y_n-A_n x_n = 0  \}.
\end{align*}
The filtering problems for the low and degenerate noise case will require sampling from  $\widetilde{\mathsf{M}}_n$ and $ \mathsf{M}_n^{\star}$ respectively. We will connect the two problems by presenting how one can view $\widetilde{\mathsf{M}}_n$  as an appropriate lifting of $ \mathsf{M}_n^{\star}$. We then discuss how this leads to generic sampling on these manifolds and then implementable PFs.
%decompose the space so that starting from the degenerate case  and use $d_x-d_y$ coordinates
%Notice both these spaces are shaped by $d_y$ constraints 
%We begin with the low noise case.   Let $x_n^{\star}$ solve $y_n=A_n x$,  $\text{ker}(A_n) = \{x\in\mathbb{R}^{d_x}:A_nx=0\}$ and $V_n$ is the $d_x\times(d_x-d_y)$ matrix whose columns are an orthonormal basis of
%$\text{ker}(A_n)$.  
We start from the degenerate case and let $x_n^{\star}\in \mathsf{M}_n^{\star}$ and let $V_n$ denote a $d_x\times(d_x-d_y)$ matrix whose columns form an orthonormal basis for the linear subspace $ \text{ker}(A_n)=\{x_n\in\mathbb{R}^{d_x}:A_nx_n=0\}$.
Thus, we have that:
%
\begin{align*}
 \mathsf{M}_n^{\star} = x_n^{\star}  + \mathscr{L}V_n,
\end{align*}
with $\mathscr{L}$ denoting linear span. The point to note here is that sampling on $\mathsf{M}_n^{\star}$ can performed by centering on  any $x_n^{\star}\in\mathsf{M}_n^{\star}$ and adding a random sample defined on $ \mathscr{L}V_n$. Due to the linear constraint from the observation, our state lives in a lower dimensional manifold, so one can proceed by appropriately sampling $z_n\in \mathbb{R}^{d_x-d_y}$ and then return the linear map $u(z_n)=x_n^{\star}  +V_n z_n$ as  a sample in $\mathsf{M}_n^{\star}$. 


We `uplift' the manifold $\mathsf{M}_n^{\star}$ to obtain  $\widetilde{\mathsf{M}}_n$ by working as follows. Consider starting from the case without noise present ($\varepsilon_n=0$) and augment $x_n^{\star}$ by setting  $(\widetilde{x}_n^{\Delta})^{\top}=((x_n^{\star})^{\top},0^{\top})$, which is a $(d_x+d_y)-$dimensional vector that obeys the defining constraint of $\widetilde{\mathsf{M}}_n$, i.e. $y_n=[A_n,\Delta^{1/2}\mathtt{I}_{d_y}]\widetilde{x}_n^{\Delta}$. 
Let $V_n^{\Delta}$ be a $(d_x+d_y)\times d_x$
matrix whose columns are an
 orthonormal basis of $\text{ker}([A_n,\Delta^{1/2}\mathtt{I}_{d_y}])$ so that we can write 
 \begin{align*}
\widetilde{\mathsf{M}}_n = \widetilde{x}_n^{\Delta,0}  + \mathscr{L}\big\{V_n^{\Delta}\big\}.
\end{align*}
We will set
\begin{align}
V_n^\Delta =
\begin{pmatrix}
V_n & \Delta^{1/2} A_n^\top \\
0 & -A_n A_n^\top
\end{pmatrix}
\label{eq:V_n_Delta}
\end{align}
and refer the reader for more details to the Appendix. From a practical point of view at each time $n$ one needs to sample a signal and observation component that lie on the manifold  $\widetilde{\mathsf{M}}_n$. The first step is to compute a point $x_n^{\star}$ solving $A_n x_n^{\star,0}=y_n$ and thus obtain  $\widetilde{x}_n^{\Delta}$ by augmenting with $0$. Then one should sample appropriately $\widetilde{z}\in\mathbb{R}^{d_x}$ and use the following mapping: 
\begin{align}
u_n^{\Delta}(\widetilde{z}) = \widetilde{x}_n^{\Delta} + V_n^{\Delta} \widetilde{z},\label{eq:u_Delta}
\end{align}
where we define  $u_n^{\Delta}(\widetilde{z})^{\top} = (u_n^{\Delta}(\widetilde{z},x)^{\top},u_n^{\Delta}(\widetilde{z},\epsilon)^{\top})$, with
$u_n^{\Delta}(\widetilde{z},x)\in\mathbb{R}^{d_x}$ and $u_n^{\Delta}(\widetilde{z},\epsilon)\in\mathbb{R}^{d_y}$, so that in the notation $u_n^{\Delta}(\widetilde{z},x)$ is meant to indicate the first $d_x$ dimensions of $u_n^{\Delta}(\widetilde{z})$ and  $u_n^{\Delta}(\widetilde{z},\epsilon)$ the last $d_y$ dimensions,  corresponding to signal and observation noise respectively.% and we note that $$y_n=[A_n,\Delta^{1/2}\mathtt{I}_{d_y}]u_n^{\Delta}(\widetilde{z})$$ as required.

In both the low noise and degenerate filtering we are primarily interested for the state variables and given this discussion it should be clear one can equivalently define the low noise filtering problem w.r.t  a random variable $\widetilde{Z}_n\in\mathbb{R}^{d_x}$.  In order to better connect the two problems  we will also use throughout the following notation convention:  $\widetilde{z}_n^{\top} = (z_n^{\top},\overline{z}_n^{\top})^{\top}$,  with $(z_n,\overline{z}_n)\in \mathbb{R}^{d_x-d_y}\times \mathbb{R}^{d_y}$, recalling that $z_n$ is the random variable of interest for the degenerate case.




\subsubsection{The Low Noise Case}\label{sec:ln_add}

PFs are defined to target a sequence of smoothing distributions, so using \eqref{eq:u_Delta} we will rewrite \eqref{eq:smoothing_M_n_low_noise} as a function of $\widetilde{z}_{1:n}$ . Note that due to the manifolds considered in this section being linear, the Gram matrices in (\ref{eq:Gram-e}), (\ref{eq:Gram}) are constant and we can write densities w.r.t.~Lebesgue measure.   Given this construction,  at any time $n\geq 1$ we can write the joint smoothing density of $\widetilde{z}_{1:n}\in\mathbb{R}^{nd_x}$ w.r.t.~$\bigotimes_{i=1}^n \lambda_{d_x}(d\widetilde{z}_i)$ as 
$$
\pi_n^{\Delta}(\widetilde{z}_{1:n}) \propto \prod_{k=1}^n p_k(u_k^{\Delta}(\widetilde{z}_k,\epsilon))
f_k(u_k^{\Delta}(\widetilde{z}_k,x)|u_{k-1}^{\Delta}(\widetilde{z}_{k-1},x))
$$
where we set $u_{0}^{\Delta}(z,x)=x_0$, for any $z\in\mathbb{R}^{d_x}$.  We are assuming the probability density is well-defined for any $\Delta> 0$, that is,  the normalization is finite.

In such a scenario,  one can use a standard particle filter,  which we now detail.  Let $\widetilde{q}_n^{\Delta}$ be a sequence of positive Markov transition densities on $\mathbb{R}^{d_x}$ and set for any $(\widetilde{z}',\widetilde{z})\in\mathbb{R}^{2d_x}$
$$
w_n^{\Delta}(\widetilde{z}',\widetilde{z}) = \frac{p_n(u_n^{\Delta}(\widetilde{z},\epsilon))
f_n(u_n^{\Delta}(\widetilde{z},x)|u_{n-1}^{\Delta}(z',x))}{\widetilde{q}_n^{\Delta}(\widetilde{z}|\widetilde{z}')}.
$$
As for the natural parameterization case, note that $\widetilde{q}_n^{\Delta}$ can depend on $\Delta$ and when $n=1$ we may (or may not) condition on $x_0$
in $\widetilde{q}_1^{\Delta}$ but write $\widetilde{q}_1^{\Delta}(\cdot|x_0)$ from herein.
Let $N\in\mathbb{N}$ be given and denote for any $n\geq 1$,  $\widetilde{z}_n^{1:N}=(\widetilde{z}_n^{1}\dots,\widetilde{z}_n^{N})$ which represent the particles of the PF.   It will prove useful to use the ancestor representation of the resampling indices so we will also be interested in $(\widetilde{a}_n^1,\dots,\widetilde{a}_n^N)\in\{1,\dots,N\}^N$ and again we use the convention that $\widetilde{a}_0^i=i$ for each $i\in\{1,\dots,N\}$.  We will also write $\widetilde{z}_{0}^{\widetilde{a}_{0}^j}=x_0$.  

The PF is then the following algorithm:
\begin{itemize}
\item{Initialize: For $i\in\{1,\dots,N\}$,  sample $\widetilde{Z}_1^i|x_0$ from $\widetilde{q}_1^{\Delta}(\cdot|x_0)$. Set $n=2$.}
\item{Iterate: For $i\in\{1,\dots,N\}$,  sample $\widetilde{a}_{n-1}^i\sim\text{Cat}(\cdot;N,\mathbf{w}^{\Delta}_{n-1})$ with 
$$\mathbf{w}^{\Delta}_{n-1}=\left( \frac{w_{n-1}^{\Delta}(\widetilde{z}_{n-2}^{\widetilde{a}_{n-2}^1},\widetilde{z}_{n-1}^{1})}
{\sum_{s=1}^N w_{n-1}^{\Delta}(\widetilde{z}_{n-2}^{\widetilde{a}_{n-2}^s},\widetilde{z}_{n-1}^s)},\ldots,\frac{w_{n-1}^{\Delta}(\widetilde{z}_{n-2}^{\widetilde{a}_{n-2}^N},\widetilde{z}_{n-1}^{N})}
{\sum_{s=1}^N w_{n-1}^{\Delta}(\widetilde{z}_{n-2}^{\widetilde{a}_{n-2}^s},\widetilde{z}_{n-1}^s)} \right).$$
%%% old version for checking
%$\widetilde{a}_{n-1}^i|\widetilde{z}_{n-2}^{1:N},\widetilde{z}_{n-1}^{1:N}$ i.i.d.~from
%$$
%\frac{w_{n-1}^{\Delta}(\widetilde{z}_{n-2}^{\widetilde{a}_{n-2}^i},\widetilde{z}_{n-1}^{\widetilde{a}_{n-1}^i})}
%{\sum_{s=1}^N w_{n-1}^{\Delta}(\widetilde{z}_{n-2}^{\widetilde{a}_{n-2}^s},\widetilde{z}_{n-1}^s)}.
%$$
For $i\in\{1,\dots,N\}$,  sample $\widetilde{Z}_n^i|\widetilde{z}_{n-1}^{\widetilde{a}_{n-1}^i}$ independently from $\widetilde{q}_n^{\Delta}(\cdot|\widetilde{z}_{n-1}^{\widetilde{a}_{n-1}^i})$.
Set $n=n+1$ and return to the start of the bullet point.}
\end{itemize}
Let $\varphi:\mathbb{R}^{d_x}\rightarrow\mathbb{R}$ be a bounded and measurable function (write such functions as $\mathtt{B}_b(\mathbb{R}^{d_x})$) then we can estimate:
$$
\pi_n^{\Delta}(\varphi) := \int_{\mathbb{R}^{nd_x}}\varphi(\widetilde{z}_n) \pi_n^{\Delta}(\widetilde{z}_{1:n})
\bigotimes_{i=1}^n \lambda_{d_x}(d\widetilde{z}_i)
$$
by using the estimator, obtained just before the end of the iterate step
$$
\pi_n^{\Delta,N}(\varphi) := \frac{\sum_{i=1}^N \varphi(\widetilde{z}_n^i)
w_{n}^{\Delta}(\widetilde{z}_{n-1}^{\widetilde{a}_{n-1}^i},\widetilde{z}_{n}^{i})}
{\sum_{i=1}^N
w_{n}^{\Delta}(\widetilde{z}_{n-1}^{\widetilde{a}_{n-1}^i},\widetilde{z}_{n}^{i})}.
$$
The convergence of this estimator and the choice of proposal has been extensively discussed in the literature; see for instance \cite{bain,delmoral} and the references therein.  We remark that the optimal,  in terms of minimizing the variance of the weights conditional on $\widetilde{z}_{n-1}$, is
$$
\widetilde{q}_n^{\Delta,\text{opt}}(\widetilde{z}_n|\widetilde{z}_{n-1}) = 
\frac{p_n(u_n^{\Delta}(\widetilde{z}_n,\epsilon))
f_n(u_n^{\Delta}(\widetilde{z}_n,x)|u_{n-1}^{\Delta}(\widetilde{z}_{n-1},x))}
{\int_{\mathbb{R}^{d_x}}
p_n(u_n^{\Delta}(\widetilde{z}_n,\epsilon))
f_n(u_n^{\Delta}(\widetilde{z}_n,x)|u_{n-1}^{\Delta}(\widetilde{z}_{n-1},x))\lambda_{d_x}(d\widetilde{z}_n)
}.
$$
This is tractable only in certain modelling scenarios, one of which we investigate below.  We remark that in the case that one
is able to compute the optimal proposal,  due to the one-to-one mappings from the natural parameterization to the noise parameterization to the map back to the Lebesgue measure in the linear constraint scenario,  the optimal proposal from the natural parameterization will provide an equivalent formulation.  One of the main points of the low noise method however,  is that it will permit its application in cases where the optimal proposal is not available and will be robust to low values of the noise parameter.

\subsubsection{The Degenerate Noise Case}\label{sec:dn_add}

The degenerate case can be handled in a similar manner,  which we now describe.
For any $(n,z)\in\mathbb{N}\times\mathbb{R}^{d_x-d_y}$ set
$$
u_n^{\star}(z) = x_n^{\star} + V_n z
$$
%where $V_n$ is the $d_x\times(d_x-d_y)$ matrix whose columns form an orthonormal basis of
%$\text{ker}(A_n)=\{x\in\mathbb{R}_{d_x}:A_nx = 0\}$ and
with the convention that $u_0^{\star}(z)=x_0$ for any $z\in \mathbb{R}^{d_x-d_y}$.
At any time $n\geq 1$ we can write the joint smoothing density of $z_{1:n}\in\mathbb{R}^{n(d_x-d_y)}$ w.r.t.~$\bigotimes_{i=1}^n \lambda_{d_x-d_y}(dz_i)$ as 
$$
\pi_n^{\star}(z_{1:n}) \propto \prod_{k=1}^n 
f_k(u_k^{\star}(z_k)|u_{k-1}^{\star}(z_{k-1})),
$$
where we note that $u_{0}^{\star}(z_{0}) = x_0$. Let $q_n^{\star}$ be a sequence of positive Markov transition densities on $\mathbb{R}^{d_x-d_y}$ and set for any $(z',z)\in\mathbb{R}^{2(d_x-d_y)}$
$$
w_n^{\star}(z',z) = \frac{
f_n(u_n^{\star}(z)|u_{n-1}^{\star}(z))}{q_n^{\star}(z|z')}.
$$
When $n=1$ we may (or may not) condition on (some part of) $x_0$
in $q_n^{\star}$ but again write $q_n^{\star}(\cdot|x_0)$ from herein.  Denote for any $n\geq 1$,  $z_n^{1:N}=(z_n^{1}\dots,z_n^{N})$, the particles of the PF.   The ancestor variables are $(a_n^1,\dots,a_n^N)\in\{1,\dots,N\}^N$ and we use the convention that $a_0^i=i$ for each $i\in\{1,\dots,N\}$.  We will also write $z_{0}^{a_{0}^j}=x_0$.

The PF is as follows.
\begin{itemize}
\item{Initialize: For $i\in\{1,\dots,N\}$,  sample $Z_1^i|x_0$ i.i.d.~from $q_1^{\star}(\cdot|x_0)$. Set $n=2$.}
\item{Iterate: For $i\in\{1,\dots,N\}$,  sample $a_{n-1}^i\sim\text{Cat}(\cdot;N,\mathbf{w}^\star_{n-1})$ with
$$
\mathbf{w}^\star_{n-1}=\left ( \frac{w_{n-1}^{\star}(z_{n-2}^{a_{n-2}^1},z_{n-1}^{1})}
{\sum_{s=1}^N w_{n-1}^{\star}(z_{n-2}^{a_{n-2}^s},z_{n-1}^s)},\ldots, \frac{w_{n-1}^{\star}(z_{n-2}^{a_{n-2}^N},z_{n-1}^{N})}
{\sum_{s=1}^N w_{n-1}^{\star}(z_{n-2}^{a_{n-2}^s},z_{n-1}^s)}\right).
$$
%
%$a_{n-1}^i|z_{n-2}^{1:N},z_{n-1}^{1:N}$ i.i.d.~from
%$$
%\frac{w_{n-1}^{\star}(z_{n-2}^{a_{n-2}^i},z_{n-1}^{a_{n-1}^i})}
%{\sum_{s=1}^N w_{n-1}^{\star}(z_{n-2}^{a_{n-2}^s},z_{n-1}^s)}.
%$$
For $i\in\{1,\dots,N\}$,  sample $Z_n^i|z_{n-1}^{a_{n-1}^i}$ independently from $q_n^{\star}(\cdot|z_{n-1}^{a_{n-1}^i})$.
Set $n=n+1$ and return to the start of the bullet point.}
\end{itemize}
Let $\varphi\in\mathtt{B}_b(\mathbb{R}^{d_x-d_y})$ then we can estimate:
$$
\pi_n^{\star}(\varphi) := \int_{\mathbb{R}^{nd_x}}\varphi(z_n) \pi_n^{\star}(z_{1:n})
\bigotimes_{i=1}^n \lambda_{d_x-d_y}(dz_i)
$$
by using the estimator, obtained just before the end of the iterate step
$$
\pi_n^{\star,N}(\varphi) := \frac{\sum_{i=1}^N \varphi(z_n^i)
w_{n}^{\star}(z_{n-1}^{a_{n-1}^i},z_{n}^{i})}
{\sum_{i=1}^N
w_{n}^{\star}(z_{n-1}^{a_{n-1}^i},z_{n}^{i})}.
$$
As in the low noise case,  the optimal proposal here would be
$$
q_n^{\star,\text{opt}}(z_n|z_{n-1}) = 
\frac{f_n(u_n^{\star}(z_n)|u_{n-1}^{\star}(z_{n-1}))}
{\int_{\mathbb{R}^{d_x-d_y}}
f_n(u_n^{\star}(z_n)|u_{n-1}^{\star}(z_{n-1}))\lambda_{d_x-d_y}(dz_n)
}.
$$

\subsubsection{A Low Noise Convergence Result}

We now consider the convergence,  in a sense to be made precise below,  of the PF in the low noise case to that of the degenerate noise case.  
Recall that we use the notation for $\widetilde{z}\in\mathbb{R}^{d_x}$ as $\widetilde{z}^{\top}=(z^{\top},\overline{z}^{\top})$, with $z\in\mathbb{R}^{d_x-d_y}$.
We make the following assumption.
\begin{hypA}\label{ass:1}
For any $(n,\widetilde{z},\widetilde{z}')\in\mathbb{N}\times\mathbb{R}^{2d_x}$ we have
$$
\lim_{\Delta\downarrow 0} \widetilde{q}_n^{\Delta}(\widetilde{z}|\widetilde{z}') = 
p_n(\overline{z}) q_n^{\star}(z|z').
$$
\end{hypA}
This assumption is not overly restrictive as one chooses the proposal densities.  The assumption means that as the noise disappears,  one should propose the $d_x-d_y$ random variable,  just as in the degenerate noise case and the noise variable according to the noise density.

We denote by $\mathbb{E}^{\Delta}$ as expectations w.r.t.~the law of the PF in the low noise case and 
by $\mathbb{E}^{\star}$ as expectations w.r.t.~the law of the PF in the degenerate noise case.  We can now state our main result whose proof is in Appendix \ref{app:res2}.
\begin{prop}\label{res:2}
Assume (A\ref{ass:1}).  Then for any $(n,r,\varphi)\in\mathbb{N}^2\times\mathtt{B}_b(\mathbb{R}^{d_x-d_y})$ we have
$$
\lim_{\Delta \downarrow 0}\mathbb{E}^{\Delta}\left[\left|\pi_n^{\Delta,N}(\varphi)-\pi_n^{\Delta}(\varphi)\right|^r\right]^{1/r} = 
\mathbb{E}^{\star}\left[\left|\pi_n^{\star,N}(\varphi)-\pi_n^{\star}(\varphi)\right|^r\right]^{1/r}.
$$
\end{prop}

\begin{rem}\label{rem:1}
The main purpose of Proposition \ref{res:2} is to verify that the properties of the low noise case PF are inherited,  as 
$\Delta \downarrow 0$ from that of the degenerate noise case.  In other words if one can design a good PF for the 
degenerate noise case,  then the low noise filtering problem can be effectively handled.  For instance if one has that
$$
\sup_{n\geq 1}\sup_{(z'z)\in\mathbb{R}^{2(d_x-d_y)}} w_n^{\star}(z',z) <+\infty
$$
and 
\begin{align*}
\inf_{n\geq 1} \inf_{(z'z)\in\mathbb{R}^{2(d_x-d_y)}} q_n(z|z') & >0  \\
\sup_{n\geq 1} \sup_{(z'z)\in\mathbb{R}^{2(d_x-d_y)}} q_n(z|z') & <+\infty
\end{align*}
then by standard results for PFs (e.g.~\cite{delmoral}) we have that
$$
\sup_{n\geq 1} \mathbb{E}^{\star}\left[\left|\pi_n^{\star,N}(\varphi)-\pi_n^{\star}(\varphi)\right|^r\right]^{1/r} \leq 
\frac{C}{N^{-\tfrac{1}{2}}}
$$
where $C<+\infty$ is a finite constant. Then for any fixed $n$ we would have
$$
\lim_{\Delta \downarrow 0}\mathbb{E}^{\Delta}\left[\left|\pi_n^{\Delta,N}(\varphi)-\pi_n^{\Delta}(\varphi)\right|^r\right]^{1/r} = 
\mathbb{E}^{\star}\left[\left|\pi_n^{\star,N}(\varphi)-\pi_n^{\star}(\varphi)\right|^r\right]^{1/r} \leq 
\frac{C}{N^{-\tfrac{1}{2}}}.
$$
\end{rem}

\subsection{Additive Gaussian Noise with Nonlinear Signal Dynamics}\label{sec:spec_add}

We consider the following state-space model for $n\in\mathbb{N}$ and $X_0\in\mathbb{R}^{d_x}$ given:
\begin{align*}
Y_n & = A_n X_n + \Delta^{1/2}\epsilon_n \\
X_n & =  f_n(X_{n-1}) + \nu_n,
\end{align*}
where $Y_n\in\mathbb{R}^{d_y}$,  $1\leq d_y < d_x$,  $A_n\in\mathbb{R}^{d_y\times d_x}$ is of full row rank, 
$\Delta\geq 0$
$f_n:\mathbb{R}^{d_x}\rightarrow\mathbb{R}^{d_x}$. Independently for each $n\in\mathbb{N}$ $\epsilon_n\sim\mathcal{N}_{d_y}(0,\Sigma_n)$ and independently $\nu_n\sim\mathcal{N}_{d_x}(0,\Omega_n)$ where, $\Sigma_n,\Omega_n$ are a sequence of real positive definite symmetric matrices. 
Here  $\mathcal{N}_d(\mathsf{m},\mathsf{K})$ denotes the $d-$dimensional Gaussian distribution with mean $\mathsf{m}$ and covariance matrix $\mathsf{K}$.% and for $d\in\mathbb{N}$ fixed,  $(\mu,\kappa)\in\mathbb{R}^{d}\times\mathbb{R}^{d\times d}$ ($\kappa$ is also positive definite and symmetric). 



\subsubsection{Optimal Proposals}

We now consider the optimal or bootstap proposals and the importance weights in the natural parameterization,
noise parameterization (low noise) and the degenerate noise case.   To aid the notations in the below statement
we partition the matrix $V_n^{\Delta}$ into its first $d_x$ rows denoted by the $d_x\times d_x$ matrix $V_n^{\Delta}(x)$ and its last $d_y$ rows denoted by the $d_y\times d_x$ matrix $V_n^{\Delta}(\epsilon)$.  We will also use the notation
$x_0^{\star}+V^{\Delta}_0(x)z=x_0$ 
$x_0^{\star}+V_0z=x_0$
and where we have not defined the objects on left hand side.
We have the following result whose proof is in Appendix \ref{app:inc}.
 
\begin{prop}\label{prop:prop}
Consider particle filters associated to the model at the start of Section \ref{sec:spec_add}.
\begin{enumerate}
\item{Natural parameterization with bootstrap proposal:
\begin{align*}
X_n | x_{n-1} & \sim \mathcal{N}_{d_x}(f(x_n),\Omega_n) \\
w_n(x_n) & = \exp\left\{-\frac{1}{2\Delta}(y_n-A_nx_n)^{\top}\Sigma_n^{-1}(y_n-A_nx_n)\right\}
\end{align*}
}
\item{Low noise with optimal proposal:
\begin{align*}
\widetilde{Z}_n|\widetilde{z}_{n-1} & \sim \mathcal{N}_{d_x}(\widetilde{\mu}_n,\widetilde{\Omega}_n) \\
w_n^{\Delta,\text{\emph{opt}}}(\widetilde{z}_{n-1}) & = \exp\left\{\frac{1}{2}
\left(\widetilde{\mu}_n^{\top}\widetilde{\Omega}_n^{-1}\widetilde{\mu}_n -
\widetilde{\mathtt{m}}_n^{\top}\Omega_n^{-1}\widetilde{\mathtt{m}}_n
\right)
\right\}
\end{align*}
where
\begin{align*}
\widetilde{\mu}_n & = \widetilde{\Omega}_nV_n^{\Delta}(x)^{\top}\Omega_n^{-1}(f_n(x_{n-1}^{\star}+V_{n-1}^{\Delta}(x)z_{n-1})-x_n^{\star})
\\
\widetilde{\Omega}_n^{-1} & =  V_n^{\Delta}(\epsilon)^{\top}\Sigma_n^{-1} V_n^{\Delta}(\epsilon) + 
 V_n^{\Delta}(x)^{\top}\Omega_n^{-1} V_n^{\Delta}(x) \\
\widetilde{\mathtt{m}}_n & = f_n(x_{n-1}^{\star}+V_{n-1}^{\Delta}(x)z_{n-1}) - x_n^{\star}.
\end{align*}
}
\item{Degenerate noise with optimal proposal:
\begin{align*}
Z_n|z_{n-1} & \sim \mathcal{N}_{d_x-d_y}(\mu_n^{\star},\Omega_n^{\star}) \\
w_n^{\star,\text{\emph{opt}}}(z_{n-1}) & = \exp\left\{\frac{1}{2}
\left((\mu_n^{\star})^{\top}(\Omega_n^{\star})^{-1}\mu_n^{\star} -
(\mathtt{m}_n^{\star})^{\top}\Omega_n^{-1}\mathtt{m}_n^{\star}
\right)
\right\}
\end{align*}
where
\begin{align*}
\mu_n^{\star} & = \Omega_n^{\star}V_n^{\top}\Omega_n^{-1}(f_n(x_{n-1}^{\star}+V_{n-1}z_{n-1})-x_n^{\star})
\\
(\Omega_n^{\star})^{-1} & =  
 V_n^{\top}\Omega_n^{-1} V_n \\
\mathtt{m}_n^{\star} & = f_n(x_{n-1}^{\star}+V_{n-1}z_{n-1}) - x_n^{\star}.
\end{align*}
}
\end{enumerate}
Finally,  we have
$$
\lim_{\Delta \downarrow 0}
w_n^{\Delta,\text{\emph{opt}}}(\widetilde{z}_{n-1}) = w_n^{\star,\text{\emph{opt}}}(z_{n-1}).
$$
\end{prop}

\begin{rem}\label{rem:rmk_lin}
The significance of Proposition \ref{prop:prop} is as follows.  $\Delta$ controls the amount of noise and the low noise case is as $\Delta\downarrow 0$.   In the case of the bootstrap PF it is clear that numerically that $w_n(x_n)$ will degenerate to zero and hence this algorithm is of little use; this is unsurprising as the proposal does not take into account any information associated to the data.  In the low and degenerate noise cases (2 \& 3 in Proposition \ref{prop:prop}) one can characterize the optimal proposal and the associated weight and, as was the case in Proposition \ref{res:2},  there is a convergence of the weights from the low noise to degenerate case with a similar conclusion to Remark \ref{rem:1}.  As we noted before the optimal proposal in the natural parameterization will be equivalent to the one in the noise parameterization. 
\end{rem}

\section{Partially Observed Diffusions}\label{sec:diff}

\subsection{Diffusion Process}

We will now consider the problem of filtering diffusion processes in the low and degenerate noise case.
We consider the following diffusion process on the filtered probability space $(\Omega,\mathscr{F},\{\mathscr{F}_t\}_{t\geq 0},\mathbb{P})$.  
\begin{equation}
dX_t = \mu(X_t)dt + \sigma(X_t)dW_t\quad\quad X_0=x_0\label{eq:diff}
\end{equation}
where $X_t\in\mathbb{R}^{d_x}$,   
$\mu:\mathbb{R}^{d_x} \rightarrow\mathbb{R}^{d_x}$, 
and $\sigma:\mathbb{R}^{d_x} \rightarrow\mathbb{R}^{d_x\times d_x}$
and $\{W_t\}_{t\geq 0}$ is a standard $d_x-$dimensional Brownian motion.
Denote $\mu^j$ as the $j^{th}-$component of $\mu$, $j\in\{1,\dots,d_x\}$
and $\sigma^{j,k}$ as the $(j,k)^{th}-$component of $\sigma$, $(j,k)\in\{1,\dots,d_x\}^2$.
As the exposition will closely follow the works in \cite{diff_bridge,beskos,non_synch,stanton} we make the following assumption which is also made in most of those papers.  Below $\mathcal{C}(\mathbb{R}^{d_x})$ are the twice continuously differentiable functions on $\mathbb{R}^{d_x}$ and $\|\cdot\|$ is the Euclidean norm.
\begin{quote} 
We have $\mu^j\in \mathcal{C}^2(\mathbb{R}^{d_x})$ and $\sigma^{j,k} \in \mathcal{C}^2(\mathbb{R}^{d_x})$, for $(j,k)\in\{1,\ldots, d\}^2$. 
In addition, $\mu$ and $\sigma$ satisfy 
\begin{itemize}
\item[(i)] {\bf uniform ellipticity}: $\Sigma(x):=\sigma(x)\sigma(x)^{\top}$ is uniformly positive definite 
over $x\in \mathbb{R}^{d_x}$;
\item[(ii)] {\bf globally Lipschitz}:
There exists a constant $C<\infty$ such that 
$|\mu^j(x)-\mu^j(x')|+|\sigma^{j,k}(x)-\sigma^{j,k}(x')| \leq C \|x-x'\|$ 
for all $(x,x') \in \mathbb{R}^{d_x}\times\mathbb{R}^{d_x}$, $(j,k)\in\{1,\dots,d_x\}^2$. 
\end{itemize}
\end{quote}
We note that one could use the Euler-Maruyama time-discretization method to obtain the transition density of the hidden process.  However,  it is well-known (e.g.~\cite{beskos} or \cite[Figure 1]{shouto}) that this method, when combined with a PF,  typically collapses (due to increasing Monte Carlo variance) as the time-discretization becomes finer and thus we consider a robust reformulation using a bridge construction that avoids this issue and is often used in the literature \cite{schauer,vd_meulen_guided_mcmc}. Both approaches will need to be discretized in time, but will opt for the latter due to its robustness to fine time discretizations. As a result the PFs that we will present in Sections \ref{sec:ln_diff} and \ref{sec:dn_diff} will be robust both to the case of low observation noise \emph{and} fine time discretizations.


\subsection{Bridge Construction}

We review the method in \cite{schauer,vd_meulen_guided_mcmc} as has been used in \cite{beskos}.
The exposition closely follows that in \cite{diff_bridge,beskos,non_synch}.
Consider the case $t\in[s_1,s_2]$, $0\leq s_1<s_2\leq T$ and let
$\mathbf{X}_{[s_1,s_2]}:=\{X_{t}\}_{t\in[s_1,s_2]}$, and $\mathbf{W}_{[s_1,s_2]}:=\{W_{t}\}_{t\in[s_1,s_2]}$.  
We will use a measure disintegration of the diffusion path measure with respect to a sampled value at time $s_2$.
Suppose $f_{t,s_{2}}(x'|x)$ denotes the unknown transition density from time $t$ to $s_2$ associated to \eqref{eq:diff}.
An equivalent formulation for the law of $\mathbf{X}_{[s_1,s_2]}$ would be to sample from $f_{s_1,s_2}$ to obtain $(x, x')\in \mathbb{R}^{2d_x}$ and then to interpolate between $x$ and $x'$ or "fill in" the path using an alternative diffusion that is conditioned to 
arrive at point $x'$ at time $s_2$. 
%We will
%we will explain that 
%we can interpolate these points by using a bridge process.
We will refer to this this diffusion process as the bridge process and its dynamics will be will be as in \eqref{eq:diff}, but instead its drift will be given by $\mu(x)+\Sigma(x)\nabla_x\log{f}_{t,s_2}(x'|x)$.
Let $\mathbb{P}_{x,x'}$ denote the law of the solution of the SDE \eqref{eq:diff}, on $[s_1,s_2]$, started at $x$ and conditioned to hit $x'$ at time $s_2$.

We introduce a user-specified auxiliary process $\{\tilde{X}_t\}_{t\in[s_1,s_2]}$ following:
%
\begin{align}
\label{eq:aux_SDE_tilde}
d \tilde X_{t} = \tilde \mu(t,\tilde X_{t})dt + \tilde \sigma(t,\tilde X_{t})dW_{t}, \quad t\in[s_1,s_2],\quad~\tilde{X}_{s_1} =x, 
\end{align}
%
where $\tilde \mu:[s_1,s_2]\times\mathbb{R}^{d_x}\rightarrow\mathbb{R}^{d_x}$ and $\tilde \sigma:\mathbb{R}^{d_x}\rightarrow\mathbb{R}^{d_x\times d_x}$ is
such that 
$$
\tilde \Sigma(s_2,x'):= \tilde \sigma(s_2,x') \tilde \sigma(s_2,x')^{\top} \equiv \Sigma(x').
$$ 
\eqref{eq:aux_SDE_tilde} is specified so that its transition density $\tilde{f}$ is available; see \cite[Section 2.2]{schauer} for the technical conditions on $\tilde \mu, \tilde \Sigma, \tilde f$. 
The main purpose of $\{\tilde X_t\}_{t\in[s_1,s_2]}$ is to sample $x'$ and use its transition density to construct another process $\{X_t^\circ\}_{t\in[s_1,s_2]}$ conditioned to hit $x'$ at $t=s_2$; which in turn will be an importance proposal for $\{X_t\}_{t\in[s_1,s_2]}$.  Set:
\begin{align}
\label{eq:aux_SDE}
d X^\circ_{t} = \mu_{s_2}^{\circ}(t,X^{\circ}_{t}; x')dt + \sigma(X^{\circ}_{t})dW_{t}, \quad t\in[s_1,s_2],\quad~X^{\circ}_{s_1} =x, 
\end{align}
where:
$$
\mu_{s_2}^{\circ}(t,x;x')=\mu(x)+\Sigma(x)\nabla_x\log\tilde{f}_{t,s_2}(x'|x),
$$ 
and denote by
 $\mathbb{P}^\circ_{x,x'}$ the probability law of the solution of (\ref{eq:aux_SDE}). %and similarly  $\tilde{\mathbb{P}}_{\theta,x,x'}$ for \eqref{eq:aux_SDE_tilde}. 
The SDE in \eqref{eq:aux_SDE} yields: 
\begin{equation}
\mathbf{W}\rightarrow C_{s_1,s_2}(x,\mathbf{W}_{[s_1,s_2]},x'),
\label{eq:map}
\end{equation} 
mapping the driving Wiener noise $\mathbf{W}$ to the solution of \eqref{eq:aux_SDE}, reparametering the problem from $\mathbf{X}$ to $(\mathbf{W},x')$.

The authors in \cite{schauer} prove that $\mathbb{P}_{x,x'}$ and $\mathbb{P}^\circ_{x,x'}$ are absolutely continuous w.r.t.~each other, with Radon-Nikodym derivative: 
\begin{equation}
\frac{d\mathbb{P}_{x,x'} }{d \mathbb{P}^\circ_{x,x'} }(\mathbf{X}_{[s_1,s_2]})=
\exp\Big\{  \int_{s_1}^{s_2} L_{s_2}(t,X_t)dt\Big\}  \frac{ \tilde{f}_{s_1,s_2}(x'|x)}{f_{s_1,s_2}(x'|x)},
\label{eq:density}
\end{equation}
where: 
\begin{align*}
L_{s_2}&(t,x):=\left(\mu(x)- \tilde{\mu}(t,x)\right)^{\top}\, \nabla_x\log\tilde{f}_{t,s_2}(x'|x)\\
&-\frac{1}{2}\textrm{Tr}\,\Big\{\,\big[ \Sigma(x)-\tilde{\Sigma}(t,x)\big] \big[ -\nabla_x^2\log\tilde{f}_{t,s_2}(x'|x)-\nabla_x\log\tilde{f}_{t,s_2}(x'|x)\nabla_x\log\tilde{f}_{t,s_2}(x'|x)^{\top} \big]\,\Big\} ,         
\end{align*}
with $\textrm{Tr}(\cdot)$ denoting the trace of a squared matrix. 

We also need to include the sampling of $x'$, so set %can write then the density of $\mathbf{X}_{[s_1,s_2]}$
$$
R_{s_1,s_2}(\mathbf{X}_{[s_1,s_2]}) := \frac{d\mathbb{P}_{x,x'} }{d \mathbb{P}^\circ_{x,x'} }(\mathbf{X}_{[s_1,s_2]})
f_{s_1,s_2}(x'|x),
%\frac{f_{s_1,s_2}(x'|x)}{\overline{f}_{s_1,s_2}(x'|x)}.
$$
where one should note that remarkably $R_{s_1,s_2}$ does not depend on the unknown transition density $f_{s_1,s_2}$.

%where and $\overline{f}_{s_1,s_2}(x'|x)$ is any positive (Lebesgue) probability density on $\mathbb{R}^{d_x}$.


%and write expectations w.r.t.:
%$$
%\overline{\mathbb{P}}(d(x_0,\dots,x_T), d(\mathbf{w}_{[0,1]},\dots,\mathbf{w}_{[T-1,T]})) = \nu(x_0)dx_0\prod_{k=1}^T \mathbb{W}(d\mathbf{w}_{[k-1,k]})\overline{f}_{k-1,k}(x_k|x_{k-1})dx_k
%$$
%as $\overline{\mathbb{E}}_{\theta}[\cdot]$,  where $\mathbb{W}$ is the law of Brownian motion.
%Set
%$$
%R_{s_1,s_2}(\mathbf{X}_{[s_1,s_2]}) := \frac{d\mathbb{P}_{x,x'} }{d \mathbb{P}^\circ_{x,x'} }(\mathbf{X}_{[s_1,s_2]})\frac{\tilde{f}_{s_1,s_2}(x'|x)}{\overline{f}_{s_1,s_2}(x'|x)}
%$$

\subsection{Bridges and Manifolds}

We will consider a linear observation process:
$$
Y_n  = A_n X_n + \Delta^{1/2}\epsilon_n 
$$
where $Y_n\in\mathbb{R}^{d_y}$,  $1\leq d_y < d_x$,  $A_n\in\mathbb{R}^{d_y\times d_x}$ is of full row rank, 
and independently for each $n\in\mathbb{N}$  $\epsilon_n\in\mathbb{R}^{d_y}$ are a sequence of independent noises with positive Lebesgue density $p_n$
and $\Delta\geq 0$. % and $\Sigma_n$ is a sequence of real positive definite symmetric matrices.
This is the same linear set-up as the previous section. Note that we are considering hidden state that evolves in continuous time $t$ that is observed at uniformly spaced discrete times and in particular when $t=n=1,2,3,\ldots$, i.e. at unit time intervals. This allows to simplify the notation and make it consistent with previous articles \cite{beskos}.  However,  this need not be the case and observations could in principle have a different amount of time in-between them,  or indeed be irregularly observed.

\subsubsection{The Low Noise Case}\label{sec:ln_diff}

Let $\mathbb{W}$ is the law of Brownian motion.
The probability measure associated to the smoother in continuous-time at time $n$ will be proportional to
\begin{equation}\label{eq:smooth_diff}
\prod_{k=1}^n\Big\{ 
R_{k-1,k}(C_{k-1,k}(u_{k-1}^{\Delta}(\widetilde{z}_{k-1},x),\mathbf{w}_{[k-1,k]},u_k^{\Delta}(\widetilde{z}_k,x)))
\mathbb{W}(d\mathbf{w}_{[k-1,k]})
p_k(u_k^{\Delta}(\widetilde{z}_k,\epsilon))\lambda_{d_x}(d\widetilde{z}_k)
\Big\}.%\times
\end{equation}
%$$
%\overline{f}_{k-1,k}(u_k^{\Delta}(\widetilde{z}_k,x)|u_{k-1}^{\Delta}(\widetilde{z}_{k-1},x))\lambda_{d_x}(\widetilde{z}_k)
%\Big\}.
%$$
In practice we cannot work with this and consider a time discretization. We first give the standard Euler-Maruyama time discretization of the solution to \eqref{eq:aux_SDE} (associated to a time interval $[k-1,k]$, $k\in\{1,\dots\}$) on a regular grid of spacing $\Xi_{l}=2^{-l}$, with starting point $x_{k-1}^{\circ}$ and ending point $x_{k}^{\circ}$.  The choice $\Xi_{l}=2^{-l}$, is again for consistency of presentation with \cite{beskos},  but other schemes can be used in practice.
We have,  for $j\in\{0,1\dots,\Xi_{l}^{-1}-2\}$:
\begin{align}\label{eq:disc_circ}
X_{k-1+(j+1)\Xi_{l}}^{\circ} & = X_{k-1+j\Xi_{l}}^{\circ} + 
\mu_{k}^{\circ}(k-1+j\Xi_{l},X^{\circ}_{k-1+j\Xi_{l}}; x_{k}^\circ)\Xi_{l} + 
\nonumber \\ &\qquad\qquad\qquad
\sigma(X^{\circ}_{k-1+j\Xi_{l}})\left[W_{k-1+(j+1)\Xi_{l}}-W_{k-1+j\Xi_{l}}\right].
\end{align}
Given $(x_{k-1}^{\circ},x_{k}^{\circ})$ and $\mathbf{W}_{[k-1,k]}^l=(W_{k-1+\Xi_{l}}-W_{k-1},\dots,W_{k-\Xi_l}-W_{k-2\Xi_{l}})$,  the recursion \eqref{eq:disc_circ} induces a discretized path 
$X_{k-1+\Xi_{l}}^{\circ},\dots,X_{k-\Xi_{l}}^{\circ}$ and we write such a path, including the starting and ending points with the notation
$$
C_{k-1,k}^l(x_{k-1}^{\circ},\mathbf{W}_{[k-1,k]}^l,x_{k}^{\circ}).
$$
We also need a discretization of the Radon-Nikodym derivative.  Consider 
$$
\mathbf{X}_{[k-1,k]}^l=(X_{k-1},X_{k-1+\Xi_{l}},\dots,X_k)
$$ 
then we set
$$
R_{k-1,k}^l(\mathbf{X}_{[k-1,k]}^l) := 
\exp\Big\{ \sum_{j=0}^{\Xi_{l}^{-1}-1} 
L_{k}(t,X_{k-1+j\Xi_{l}})\Xi_{l}\Big\}\tilde{f}_{k-1,k}(X_{k}|X_{k-1}).% \times \frac{\tilde{f}_{k-1,k}(X_{k}|X_{k-1})}{\overline{f}_{k-1,k}(X_{k}|X_{k-1})}.
$$

As a result of the time-discretization we can target the smoother that is proportional to
$$
\prod_{k=1}^n\Big\{ 
R_{k-1,k}^l(C_{k-1,k}^l(u_{k-1}^{\Delta}(\widetilde{z}_{k-1},x),\mathbf{w}_{[k-1,k]}^l,u_k^{\Delta}(\widetilde{z}_k,x)))
\mathbb{W}(d\mathbf{w}_{[k-1,k]}^l)
p_k(u_k^{\Delta}(\widetilde{z}_k,\epsilon))\lambda_{d_x}(d\widetilde{z}_k)
\Big\}%\times
$$
%$$
%\overline{f}_{k-1,k}(u_k^{\Delta}(\widetilde{z}_k,x)|u_{k-1}^{\Delta}(\widetilde{z}_{k-1},x))\lambda_{d_x}(\widetilde{z}_k)
%\Big\}.
%$$
For the importance sampling part, we proceed similar to Section \ref{sec:ln_add} and for the weights set:
\begin{align*}
w_n^{\Delta,l}(\widetilde{z}_{n-1},\mathbf{w}_{[n-1,n]}^l,\widetilde{z}_n)&  = 
R_{n-1,n}^l(C_{n-1,n}^l(u_{n-1}^{\Delta}(\widetilde{z}_{n-1},x),\mathbf{w}_{[n-1,n]}^l,u_n^{\Delta}(\widetilde{z}_n,x)))
%\times \\
%&
\frac{
p_n(u_n^{\Delta}(\widetilde{z}_n,\epsilon))}%\overline{f}_{n-1,n}(u_n^{\Delta}(\widetilde{z}_n,x)|u_{n-1}^{\Delta}(\widetilde{z}_{n-1},x))}
{\widetilde{q}_n^{\Delta}(\widetilde{z}_n|\widetilde{z}_{n-1})}.
\end{align*}
We can now modify the PF of Section \ref{sec:ln_add} to this case as follows.
\begin{itemize}
\item{Initialize: For $i\in\{1,\dots,N\}$,  sample $\widetilde{Z}_1^i|x_0$ i.i.d.~from $\widetilde{q}_1^{\Delta}(\cdot|x_0)$.
and independently $\mathbf{w}_{[0,1]}^{i,l}\sim\mathbb{W}(\cdot)$.
 Set $n=2$.}
\item{Iterate: For $i\in\{1,\dots,N\}$,  sample $\widetilde{a}_{n-1}^i\sim\text{Cat}(\cdot;N,\mathtt{w}^\Delta_n)$ with

$$\mathtt{w}^\Delta_n=\left(\frac{w_{n-1}^{\Delta,l}(\widetilde{z}_{n-2}^{\widetilde{a}_{n-2}^1}, \mathbf{w}_{[n-2,n-1]}^{\widetilde{a}_{n-2}^1,l},\widetilde{z}_{n-1}^{1})}
{\sum_{s=1}^N w_{n-1}^{\Delta,l}(\widetilde{z}_{n-2}^{\widetilde{a}_{n-2}^s},
\mathbf{w}_{[n-2,n-1]}^{\widetilde{a}_{n-2}^s,l},
\widetilde{z}_{n-1}^s)},\ldots,
\frac{w_{n-1}^{\Delta,l}(\widetilde{z}_{n-2}^{\widetilde{a}_{n-2}^N}, \mathbf{w}_{[n-2,n-1]}^{\widetilde{a}_{n-2}^N,l},\widetilde{z}_{n-1}^{N})}
{\sum_{s=1}^N w_{n-1}^{\Delta,l}(\widetilde{z}_{n-2}^{\widetilde{a}_{n-2}^s},
\mathbf{w}_{[n-2,n-1]}^{\widetilde{a}_{n-2}^s,l},
\widetilde{z}_{n-1}^s)}
\right)$$
%
%
%$\widetilde{a}_{n-1}^i|\widetilde{z}_{n-2}^{1:N},
%\mathbf{w}_{[n-2,n-1]}^{1:N,l},\widetilde{z}_{n-1}^{1:N}$ i.i.d.~from
%$$
%\frac{w_{n-1}^{\Delta,l}(\widetilde{z}_{n-2}^{\widetilde{a}_{n-2}^i}, \mathbf{w}_{[n-2,n-1]}^{\widetilde{a}_{n-2}^i,l},\widetilde{z}_{n-1}^{\widetilde{a}_{n-1}^i})}
%{\sum_{s=1}^N w_{n-1}^{\Delta,l}(\widetilde{z}_{n-2}^{\widetilde{a}_{n-2}^s},
%\mathbf{w}_{[n-2,n-1]}^{\widetilde{a}_{n-2}^s,l},
%\widetilde{z}_{n-1}^s)}.
%$$
For $i\in\{1,\dots,N\}$,  sample independently $\widetilde{Z}_n^i|\widetilde{z}_{n-1}^{\widetilde{a}_{n-1}^i}$  from $\widetilde{q}_n^{\Delta}(\cdot|\widetilde{z}_{n-1}^{\widetilde{a}_{n-1}^i})$
and $\mathbf{w}_{[n-1,n]}^{i,l}\sim\mathbb{W}(\cdot)$.
Set $n=n+1$ and return to the start of the bullet point.}
\end{itemize}


\subsubsection{Degenerate Noise}\label{sec:dn_diff}

In the degenerate noise case,  one can follow the same path as the low noise case and write the time discretized smoother
up-to a normalizing constant as
$$
\prod_{k=1}^n\Big\{ 
R_{k-1,k}^l(C_{k-1,k}^l(u_{k-1}^{\star}(z_{k-1}),\mathbf{w}_{[k-1,k]}^l,u_k^{\star}(z_k)))
\mathbb{W}(d\mathbf{w}_{[k-1,k]}^l)
%\overline{f}_{k-1,k}(u_k^{\star}(z_k)|u_{k-1}^{\star}(z_{k-1}))
\lambda_{d_x-d_y}(dz_k)
\Big\}.
$$
Similar to Section \ref{sec:dn_add}, for the importance weights we set:
\begin{align*}
w_n^{\star,l}(z_{n-1},\mathbf{w}_{[n-1,n]}^l,z_n)&  = 
%\times \\
%&
\frac{R_{n-1,n}^l(C_{n-1,n}^l(u_{n-1}^{\star}(z_{n-1}),\mathbf{w}_{[n-1,n]}^l,u_n^{\star}(z_n)))
}{q_n^{\star}(z_n|z_{n-1})}.%\overline{f}_{n-1,n}(u_n^{\star}(z_n)|u_{n-1}^{\star}(z_{n-1},x))
\end{align*}
We can now modify the PF of Section \ref{sec:dn_add} to this case as follows.
\begin{itemize}
\item{Initialize: For $i\in\{1,\dots,N\}$,  sample $Z_1^i|x_0$ i.i.d.~from $q_1^{\star}(\cdot|x_0)$ and independently $\mathbf{w}_{[0,1]}^{i,l}\sim\mathbb{W}(\cdot)$. Set $n=2$.}
\item{Iterate: For $i\in\{1,\dots,N\}$,  Sample $a_{n-1}^i\sim\text{Cat}(\cdot;N,\mathtt{w} ^{\star,l}_{n-1})$ i.i.d. with
% $a_{n-1}^i|z_{n-2}^{1:N},\mathbf{w}_{[n-2,n-1]}^{1:N,l},z_{n-1}^{1:N}$ i.i.d.~from
$$
\mathtt{w} ^{\star,l}_{n-1}=\left(
\frac{w_{n-1}^{\star,l}(z_{n-2}^{a_{n-2}^1},
\mathbf{w}_{[n-2,n-1]}^{a_{n-2}^1},z_{n-1}^1)}
{\sum_{s=1}^N w_{n-1}^{\star,l}(z_{n-2}^{a_{n-2}^s},\mathbf{w}_{[n-2,n-1]}^{a_{n-2}^s},z_{n-1}^s)},\ldots,\frac{w_{n-1}^{\star,l}(z_{n-2}^{a_{n-2}^N},
\mathbf{w}_{[n-2,n-1]}^{a_{n-2}^N},z_{n-1}^N)}
{\sum_{s=1}^N w_{n-1}^{\star,l}(z_{n-2}^{a_{n-2}^s},\mathbf{w}_{[n-2,n-1]}^{a_{n-2}^s},z_{n-1}^s)}\right).
$$
For $i\in\{1,\dots,N\}$,  sample independently $Z_n^i|z_{n-1}^{a_{n-1}^i}$  from $q_n^{\star}(\cdot|z_{n-1}^{a_{n-1}^i})$
and $\mathbf{w}_{[n-1,n]}^{i,l}\sim\mathbb{W}(\cdot)$.
Set $n=n+1$ and return to the start of the bullet point.}
\end{itemize}
The algorithm here is presumably a PF implementation based on the ideas in \cite{bierkens} for the uniformly elliptic case.


\section{Numerical Results}\label{sec:numerics}

For all numerical examples we have
\begin{equation*}
    Y_n = A_nX_n + \Delta^{1/2}\epsilon_n
\end{equation*}
where,  independently for each $n\in\mathbb{N}$,  $\epsilon_n\sim\mathcal{N}_{d_y}(0,\mathtt{I}_{d_y})$ and varying $\Delta$.   In Section \ref{ex:lin_gauss} $d_y=1$,  in Section \ref{ex:non_lin} $d_y=2$ and in Section \ref{ex:diff} $d_y=1$.

\subsection{Linear Gaussian Model}\label{ex:lin_gauss}

We consider a linear Gaussian model with $d_x=10$
\begin{equation*}
    X_n = B_nX_{n-1} + \nu_n
\end{equation*}
where,  independently for each $n\in\mathbb{N}$,  $\nu_n\sim\mathcal{N}_{d_x}(0,\Omega)$. We choose $X_0=(0,\dots,0)^\top$, $\Omega=\mathtt{I}_{d_x}$, $B_n=0.9\mathtt{I}_{d_x}$, $(A_n)_{ij}=1/d_x$ (computes the average of the state),  $n\in\{1,\dots,20\}$ and $\Delta=10^{-4}$. We run the Kalman filter, the particle filters with bootstrap and optimal proposals under the natural parameterization and the proposed filter under the noise parameterization. We set the number of particles to $N=10^4$ and resample if the effective sample size (ESS) drops below $N/2$.  In Figure \ref{fig:lgm_1}, we plot the marginal densities of the first component of $X_n$ at various $n$. We see that, except the bootstrap, all filters yield similar estimates as one might expect in this case. 

\begin{figure}[h]
    \centering
    \includegraphics[width=\linewidth]{figures/pf_lgm_1.pdf}
    \caption{Marginals densities of $X_{n,1}$ for $n\in\{1,6,11,16\}$.  This is for the linear Gaussian model.}
    \label{fig:lgm_1}
\end{figure}

We now repeat the experiment for various $\Delta$ and alllow $\Delta\to0$.  We monitor the median ESS (Figure \ref{fig:lgm_3}) across the time-steps for a given $\Delta$.
In Figure \ref{fig:lgm_3}
we observe that the bootstrap PF degenerates for relatively large $\Delta$,  whilst the other methods remain stable. %The optimal proposal under natural parameterization becomes unstable for very tiny $\Delta$,  due to numerical errors and the reasoning in Remark \ref{rem:rmk_lin}.  
We can also see it in the mean squared error (MSE) of the mean estimates (error averaged across time and dimensions, see Equation \ref{eq:mse}), where the bootstrap PF collapses but the low-noise optimal proposal converges to the degenerate case.
\begin{equation}\label{eq:mse}
    \text{MSE} = \frac{1}{n_{\text{step}}}\sum_{n=1}^{n_{\text{step}}}\left\{ \frac{1}{d_x}\sum_{j=1}^{d_x} \left[ \left( \frac{1}{N}\sum_{i=1}^{N}X_{n,j}^{i} \right)-\mathbb{E}\left( X_{n,j} \right) \right]^2 \right\}
\end{equation}

\begin{figure}[h]
    \centering
    \includegraphics[width=\linewidth]{figures/pf_lgm_3.pdf}
    \caption{ESS (left) and MSE (right) versus $\log_{10}(1/\Delta)$.  This is for the linear Gaussian model.}
    \label{fig:lgm_3}
\end{figure}


\subsection{Nonlinear Model with Additive Noise}\label{ex:non_lin}

Now we consider a nonlinear Lorenz-96 model with additive $t$-distributed noise, $d_x=8$
\begin{equation*}
    X_n = F(X_{n-1}) + \nu_n
\end{equation*}
where $F_j(x) = x_j + \Delta t ((x_{j+1}-x_{j-2})x_{j-1}-x_j+F_0)$, $j\in\{1,\dots,d_x\}$ and,  independently for each $n\in\mathbb{N}$, $\nu_n\sim t_{d_x}(0,\Omega,\vartheta)$ (multivariate $t$-distribution with scale matrix $\Omega$ and $\vartheta>0$ degrees of freedom). We fix the parameter of the model to $F_0=8$ and the degrees of freedom of the innovation noise to $\vartheta=2$, and simulate it at stepsize $\Delta t=10^{-2}$ for $n\in\{1,\dots,400\}$ to obtain the data, starting from $X_0=(1,\dots,1)^\top$.  We fix $\Delta=10^{-4}$ and $\Omega=\Delta t\mathtt{I}_{d_x}$, the matrix $A_n$ is constructed such that we observe the $1^{\text st}$ and $5^{\text th}$ component of $X_n$. We can observe the chaotic behavior of the model in Figure \ref{fig:lorenz_0},  which are the simulated dynamics of the model.

\begin{figure}[h]
    \centering
    \includegraphics[width=\linewidth]{figures/pf_lorenz_0.pdf}
    \caption{Components of $X_n$ for $n\in\{1,\dots,400\}$. This is for the nonlinear model.}
    \label{fig:lorenz_0}
\end{figure}

Note that the optimal proposal is not available directly due to $\nu_n$ being $t$-distributed. However, it admits a hierarchical representation as a scale-mixture of Gaussians, via $\nu_n|\omega_n\sim\mathcal{N}(0,\Omega/\omega_n)$ with a Gamma-distributed auxiliary variable $\omega_n\sim\Gamma(\vartheta/2,\vartheta/2)$. Given $\omega_n$ the noise is conditionally Gaussian, so can use the optimal proposals from the Gaussian case to obtain a \emph{conditionally} optimal proposal for both the natural and noise parameterizations, conditioned on $\omega_n$, which we sample from the prior $\Gamma(\vartheta/2,\vartheta/2)$. 

We run all the mentioned filters (except for Kalman filter which does not apply here) again for $N=10^4$ particles with the same resampling rule.  We plot several marginals in Figure \ref{fig:lorenz_1}, where we recover similar results as in the case of the linear Gaussian model, except for the case of the bootstrap PF whose performance is too poor to give presentable results. Further, we plot the ESS versus the time index in Figure \ref{fig:lorenz_2}, where we see that the bootstrap PF  collapses from the very beginning, while other methods give adequate results in terms of this measure.
To further examine the accuracy of the PFs we plot individual sample paths of the particles (for a given run) and their mean for $8^{\text th}$ component in Figure \ref{fig:lorenz_3}.  In both the top and bottom of  Figure \ref{fig:lorenz_3} the true value is the actual signal used to generate the data.   Apart from the bootstrap PF all of the methods are seemingly able to recover the true hidden state in the $8^{\text th}$ component.
In Figure \ref{fig:lorenz_5} we consider re-running all of the methods for different $\Delta$ and sending $\Delta\to 0$. 
Figure \ref{fig:lorenz_5} displays the median ESS (across time) and the MSE (between estimated mean and the true signal) of the four PFs.
In this case the optimal proposal under the natural parameterization does not seem to collapse (in terms of ESS)  but this could be as the weight does not degenerate to zero. The MSE of the bootstrap is too large as it collapses immediately. The other two seem to roughly approach the MSE of the degenerate PF, but there is still variation due to iterates not being exactly the same.

\begin{figure}[h]
    \centering
    \includegraphics[width=\linewidth]{figures/pf_lorenz_1.pdf}
    \caption{Marginals densities of $X_{n,8}$ for $n\in\{1,101,201,301\}$. This is for the nonlinear model.}
    \label{fig:lorenz_1}
\end{figure}

\begin{figure}[h]
    \centering
    \includegraphics[width=0.6\linewidth]{figures/pf_lorenz_2.pdf}
    \caption{ESS versus $n$.  This is for the nonlinear model.}
    \label{fig:lorenz_2}
\end{figure}


\begin{figure}[h]
    \centering
    \includegraphics[width=\linewidth]{figures/pf_lorenz_3.pdf}
    \includegraphics[width=\linewidth]{figures/pf_lorenz_4.pdf}
    \caption{Sample paths (top) and mean of the filter (bottom) for $X_{n,8}$.  This is for the nonlinear model.}
    \label{fig:lorenz_3}
\end{figure}


\begin{figure}[h]
    \centering
    \includegraphics[width=\linewidth]{figures/pf_lorenz_5.pdf}
    \caption{ESS (left) and MSE (right) versus $\log_{10}(1/\Delta)$. This is for the nonlinear model.}
    \label{fig:lorenz_5}
\end{figure}


\subsection{Diffusion Model}\label{ex:diff}

We now consider the FitzHugh-Nagumo diffusion model with $d_x=2$:
\begin{equation*}
    dX_t = \mu(X_t)dt + \sigma(X_t)dW_t   
\end{equation*}
with $\mu(x)=((x_1-x_1^3-x_2)/\alpha, \gamma x_1-x_2+\beta)^\top$ and $\sigma(x)=\sigma_0 \mathtt{I}_{d_x}$. This is an elliptic diffusion with constant diffusion term.  We fix the model parameters to $\alpha=0.1$, $\gamma=1$, $\beta=0.2$ and $\sigma_0=0.1$. We choose the starting values to be $X_0=(0.5, 0.5)^\top$. We simulate the dynamics in the interval $t\in[0,10]$ using the Euler-Maruyama scheme with stepsize $\Xi=5\times10^{-3}$ and observe the first component at regular times with spacing $10^{-2}$ and finally $\Delta=10^{-8}$ is the noise level. This scenario leads to 20 discretization points between two consecutive observations and 2000 in total. The simulated dynamics is given in Figure \ref{fig:fhn_0}.

\begin{figure}[h]
    \centering
    \includegraphics[width=\linewidth]{figures/pf_fhn_0.pdf}
    \caption{Components of $X_t$ for $t\in[0,10]$. This is for the diffusion model.}
    \label{fig:fhn_0}
\end{figure}

We consider an auxiliary process with linearized drift $\tilde{\mu}(x)=\mu(x_{\text{ref}})+\frac{\partial}{\partial x}\mu(x_{\text{ref}})(x-x_{\text{ref}})$.  We can calculate the reference point $x_{\text{ref}}$ by evolving the true dynamics deterministically starting from $x$ using the stepsize $\Xi$ for 20 iterations until we reach the next observation time. 
The auxiliary process is an Ornstein-Uhlenbeck process with a transition density of $x'$ given $x$ over time step $h$ as $\mathcal{N}_2(m, Q)$ where
\begin{equation}\label{eq:diff_aux_transition}
\begin{aligned}
    & m=\exp(Jh)x+J^{-1}(\exp(Jh)-I)(\mu-Jx_{\text{ref}}) \\
    & Q=\exp(Jh)\Psi\exp(J^\top h)-\Psi
\end{aligned}
\end{equation}
where $J=\frac{\partial}{\partial x}\mu(x_{\text{ref}})$, $\mu=\mu(x_{\text{ref}})$ and $\text{vec}(\Psi)=(I\otimes J+J\otimes I)^{-1}\text{vec}(\Sigma)$, $\otimes$ is a Kronecker product and $\exp(Jh)$ is a matrix exponentiation. Here, we assume that the jacobian $J$ is both diagonalizable and invertible.

We consider 4 PFs, the low noise and degenerate cases have been detailed. We also consider a `bootstrap' and `optimal' PF which are a little different from the conventional ones we have used so far.  Both the bootstrap and
optimal methods consider a smoother of the type \eqref{eq:smooth_diff} except under the natural parameterization.
That is, after time discretization,  the target smoother is up-to a constant at time $n$
$$
\prod_{k=1}^n\Big\{ h_k(y_k|x_k)
R_{k-1,k}^l(C_{k-1,k}^l(x_{k-1},\mathbf{w}_{[k-1,k]}^l,
x_k))
\mathbb{W}(d\mathbf{w}_{[k-1,k]}^l)
\lambda_{d_x}(x_k)
\Big\}.
$$
The bootstrap PF will use a proposal that is exactly the auxiliary process (on $x_k$ the Brownian motion is sampled from its law)  whereas the optimal (which is not optimal in any sense) is a proposal proportional to the term $h_k(y_k|x_k)$ multiplied by the density of the auxiliary process given by Equation \ref{eq:diff_aux_transition} and the Brownian motion sampled from its law.


In Figure \ref{fig:fhn_2} we plot the ESS against time for a particular run. Here, we have set $N=10^4$ and resample at each step.  As in the previous example, the bootstrap PF collapses completely and the other filters perform quite well for this criterion.  This shows that one can filter in low or degenerate noise and with a very fine time discretization.  This deduction is supported by the results in Figures \ref{fig:fhn_1} (marginal densities of the state) and \ref{fig:fhn_3} (tracking the true signal) which shows that one can reconstruct the signal process in this case. In Figure \ref{fig:fhn_5} we repeat the experiment for $\Delta\to0$, where we see how the bootstrap PF collapses and the low-noise PF converges to the degenerate case.

\begin{figure}[h]
    \centering
    \includegraphics[width=0.6\linewidth]{figures/pf_fhn_2.pdf}
    \caption{ESS versus $n$.  This is for the diffusion model.}
    \label{fig:fhn_2}
\end{figure}

\begin{figure}[h]
    \centering
    \includegraphics[width=\linewidth]{figures/pf_fhn_1.pdf}
    \caption{Marginals of $X_{t,2}$ for $n\in\{1, 34, 67, 100\}$.  This is for the diffusion model.}
    \label{fig:fhn_1}
\end{figure}

\begin{figure}[h]
    \centering
    \includegraphics[width=\linewidth]{figures/pf_fhn_3.pdf}
    \includegraphics[width=\linewidth]{figures/pf_fhn_4.pdf}
    \caption{Sample paths (top) and the mean estimates (bottom) versus time. This is for the diffusion model.}
    \label{fig:fhn_3}
\end{figure}

\begin{figure}[h]
    \centering
    \includegraphics[width=\linewidth]{figures/pf_fhn_5.pdf}
    \caption{ESS (left) and MSE (right) versus $\log_{10}(1/\Delta)$. This is for the diffusion model.}
    \label{fig:fhn_5}
\end{figure}

\subsubsection{State-dependent Diffusion Model}\label{ex:diff_2}

A more interesting example is when the diffusion term in FitzHugh-Nagumo model depends on the state.  We now revisit the previous model with exactly the same setting, except with the diffusion coefficient
\begin{equation*}
    \sigma(x) = \begin{bmatrix}
        \sigma_1\sqrt{x_1^2+1} & 0 \\  0 & \sigma_2\sqrt{x_2^2+1}
    \end{bmatrix}
\end{equation*}
where we fix $\sigma_1=\sigma_2=0.1$.  We remark that,  contrary to the previous example (constant diffusion coefficient),  we do not know of any potential alternative PF method.  In the constant diffusion case,  one could envisage using the Girsanov method w.r.t.~a linear diffusion with the same constant drift,  but that is not possible in this case.  The auxiliary process is as before except with $\tilde{\sigma}=\sigma(x')$, $x'$ the end-point. 
To actually propose the new position of the particle instead of using the auxiliary process, 
we use an approximation of the auxiliary density by substituting $x'$ with deterministic $x_{\text{ref}}$ and using $\sigma(x_{\text{ref}})$.   This incurs an extra correction term in the importance weights as now one needs to divide by the associated density.  The `bootstrap' and `optimal' methods use this new density in lieu of the auxiliary process density used in the constant diffusion case.  We run each PF with $N=10^4$ particles.

We consider the ESS versus $\log_{10}(1/\Delta)$ in Figure \ref{fig:fhndiff_5}.  We observe similar results to the case of a constant diffusion coefficient.   This shows that one can use the method for non-constant diffusions in low noise and high time discretization error.

\begin{figure}[h]
    \centering
    \includegraphics[width=\linewidth]{figures/pf_fhndiff_5.pdf}
    \caption{ESS (left) and MSE (right) versus $\log_{10}(1/\Delta)$. This is for the state-dependent diffusion model.}
    \label{fig:fhndiff_5}
\end{figure}

\section{Conclusions}

We have presented a novel PF framework that can be very effective for likelihoods with low and even degenerate observational noise. We have focused on the linear observation case and our numerical results showed our method is very accurate in a variety of settings. Our method takes advantage of being able to compute $V_n$ or $V_n^\Delta$ and having an explicit maps $u_k^\star$ or $u_k^\Delta$, so extending this PF methodology for the general non-linear manifolds whilst challenging would be interesting future work. Finally, the presented reparametrizations should be also useful for Sequential Monte Carlo methods applied to Approximate Bayesian Computation and interpreting $\Delta$ as a a tolerance threshold. 

\appendix

\section{Proofs and Auxiliary Material}\label{app:proof}

\subsection{On construction of $V_n^\Delta$}

We will present motivation for the specific choice for $V_n^{\Delta}$ in \eqref{eq:V_n_Delta}. When choosing a possible basis for $\text{ker}([A_n,\Delta^{1/2}\mathtt{I}_{d_y}])$, a common choice would be to choose $\{\left(u_i^\top,(-{\Delta}^{-1/2}Au_i)^\top\right)^\top\}_{i=1}^{d_x}$ with $\Delta>0$ and $\{u_i\}_{i=1}^{d_x}$ denoting a standard basis for $\mathbb{R}^{d_x}$. For $\Delta=0$, it is clear that the last $d_y$ coordinates due to $\epsilon$ should be linearly independent as they do not enter the constraint, i.e. $\text{ker}([A_n,\mathtt{0}_{d_y}])=\text{ker}(A_n)\times\mathbb{R}^{d_y}$. This is not convenient for us, as we would like to avoid having to consider separate cases for $\Delta\neq0$ and $\Delta=0$ and more importantly we would like to express more directly $V_n^\Delta$ in terms of $V_n$. For this we first note that $\mathbb{R}^{d_x}=\text{ker}(A)\oplus\text{range}(A^{\top})$ and as $V_n$ spans $\text{ker}(A)$ one can use the row span of $A$ instead  for the space of the coordinates related to $\epsilon$ . Henceforth denote the row vectors of $A_n$ as  $a_{n,1},\ldots, a_{n,{d_y}}$ and recall they are assumed linearly independent with $A$ having a column rank being $d_y$. With this decomposition of one can identify how to augment the space. If $x\in\text{ker}(A)$ we need to augment with $0$ and otherwise use $(x,\epsilon)=(\Delta^{1/2}x,-Ax)$. Let $V_n=[ e_{n,1},\ldots, e_{n,d_x-d_y} ]\in \mathbb{R}^{d_x\times(d_x-d_y)}$ whose columns $\{e_{n,i}\}_{i=1}^{d_x-d_y}$ are an orthonormal basis for the linear subspace $ \text{ker}(A_n)$. Then define the following two {orthogonal} sets of vectors: 
%
\begin{align*}
\widetilde{E}_n^{\Delta} &=\big\{ (e_{n,1}^{\top},0^{\top})^{\top},\ldots,   (e_{n,d_x-d_y}^{\top},0^{\top})^{\top}\big\};
 \\
\widetilde{F}_n^{\Delta} & = \big\{ (\Delta^{1/2}a_{n,1}, (-A_n a_{n,1}^{\top})^{\top})^{\top},\ldots, (\Delta^{1/2}a_{n,d_y}, (-A_n a_{n,d_y}^{\top})^{\top})^{\top}\big\}. 
\end{align*}
As a result we end up with the expression in \eqref{eq:V_n_Delta}, that we recall here
\begin{align*}
V_n^\Delta =
\begin{pmatrix}
V_n & \Delta^{1/2} A_n^\top \\
0 & -A_n A_n^\top
\end{pmatrix}.
\end{align*}
%%As a result we have: 
%\begin{align*}
%\widetilde{\mathsf{M}}_n^{\Delta} = \widetilde{x}_n^{\Delta}  + \mathscr{L}\big\{\widetilde{E}_n^{\Delta}\cup \widetilde{F}_n^{\Delta} \big\} , \qquad \widetilde{E}_n^{\Delta} \perp  \widetilde{F}_n^{\Delta},
%\end{align*}
%%i.e.  $V_n^{\Delta}\in \mathbb{R}^{(d_x+dy)\times d_x}$ a matrix with columns made of the elements of  $\widetilde{E}_n^{\Delta}\cup \widetilde{F}_n^{\Delta}$. 
Note that when $\Delta=0$ this is consistent with $\text{ker}([A_n,\Delta^{1/2}\mathtt{I}_{d_y}])=\text{ker}(A_n)\times\mathbb{R}^{d_y}$.
%, where one could equivalently write instead
%\begin{align*}
%V_n^\Delta =
%\begin{pmatrix}
%V_n & \Delta^{1/2} A_n^\top \\
%0 & -I_{d_y}^\top
%\end{pmatrix}.
%\end{align*}


\subsection{Proof of Proposition \ref{res:2}}\label{app:res2}

%Recall ,  we set $\widetilde{z}_n^{\top} = (z_n^{\top},\overline{z}_n^{\top})^{\top}$,  with $(z_n,\overline{z}_n)\in \mathbb{R}^{d_x-d_y}\times \mathbb{R}^{d_y}$.

\begin{proof}
Suppose $\widetilde{z}_n^{\top} = (z_n^{\top},\check{z}_n^{\top})^{\top}$,  with $(z_n,\check{z}_n)\in \mathbb{R}^{d_x-d_y}\times \mathbb{R}^{d_y}$. Based on the above discussion  we have that for any fixed $\widetilde{z}_n\in\mathbb{R}^{d_x}$ %mentioned above was recognized in \cite{zhu}
\begin{equation}\label{eq:prf1}
u_n^{\Delta}(\widetilde{z}) \rightarrow
\begin{bmatrix}
x_n^{\star} \\
0
\end{bmatrix} + 
\begin{bmatrix}
V_n z_n  \\
\bar{z}_n
\end{bmatrix},\quad\text{as}\quad{\Delta\downarrow 0},
\end{equation}
where we use the linear transformation $\overline{z}_n=-A_nA_n^\top \check
{z}_n$.
We remark that for any $\varphi\in\mathtt{B}_b(\mathbb{R}^{d_x-d_y})$
$$
\int_{\mathbb{R}^{nd_x}} \varphi(z_n)  \prod_{k=1}^n p_k(u_k^{\Delta}(\widetilde{z}_k,\epsilon))
f_k(u_k^{\Delta}(\widetilde{z}_k,x)|u_{k-1}^{\Delta}(\widetilde{z}_{k-1},x)) 
\bigotimes_{i=1}^n \lambda_{d_x}(d\widetilde{z}_i)
$$
integrates to a finite constant that is upper-bounded uniformly $\Delta$ as $p_k$ and $f_k$ are probability densities,
so one can change variables and the resulting jacobians are upper-bounded in $\Delta$.
Therefore,  by the dominated convergence theorem we have for any $\varphi\in\mathtt{B}_b(\mathbb{R}^{d_x-d_y})$
$$
\lim_{\Delta\downarrow 0} \int_{\mathbb{R}^{nd_x}} \varphi(z_n)  \prod_{k=1}^n p_k(u_k^{\Delta}(\widetilde{z}_k,\epsilon))
f_k(u_k^{\Delta}(\widetilde{z}_k,x)|u_{k-1}^{\Delta}(\widetilde{z}_{k-1},x)) 
\bigotimes_{i=1}^n \lambda_{d_x}(d\widetilde{z}_i)
 = 
$$
$$
\int_{\mathbb{R}^{nd_x}} \varphi(z_n)  \prod_{k=1}^n p_k(\overline{z}_k)
f_k(u_k^{\star}(z_k)|u_{k-1}^{\star}(z_{k-1})) 
\bigotimes_{i=1}^n \lambda_{d_x}(d\widetilde{z}_i) 
 =
$$
$$ 
 \int_{\mathbb{R}^{n(d_x-d_y)}}
\varphi(z_n)  \prod_{k=1}^n 
f_k(u_k^{\star}(z_k)|u_{k-1}^{\star}(z_{k-1})) 
\bigotimes_{i=1}^n \lambda_{d_x}(dz_i)
$$
and hence that
$$
\lim_{\Delta\downarrow 0}\pi_n^{\Delta}(\varphi) = \pi_n^{\star}(\varphi).
$$

Now,  using \eqref{eq:prf1} and assumption (A\ref{ass:1}) we have for any fixed $(\widetilde{z}',\widetilde{z})\in\mathbb{R}^{d_x}$
\begin{equation}\label{eq:prf2}
\lim_{\Delta\downarrow 0} w_n^{\Delta}(\widetilde{z}',\widetilde{z}) = 
\frac{p_n(\overline{z})f_n(u_n^{\star}(z)|u_{n-1}^{\star}(z'))}
{p_n(\overline{z})q_n(z|z')} = w_n^{\star}(z',z).
\end{equation}
Now the joint density of the PF in the low noise case is
$$
\mathbb{P}^{\Delta}_n(\widetilde{z}_{1:n}^{1:N},\widetilde{a}_{1:n-1}^{1:N})
= \left\{\prod_{i=1}^N \widetilde{q}_1^{\Delta}(\widetilde{z}_1^i|x_0)\right\}\left\{
\prod_{k=2}^n \prod_{i=1}^N \frac{w_{k-1}^{\Delta}(\widetilde{z}_{k-2}^{\widetilde{a}_{k-2}^i},\widetilde{z}_{k-1}^{\widetilde{a}_{k-1}^i})}{
\sum_{s=1}^N
w_{k-1}^{\Delta}(\widetilde{z}_{k-2}^{\widetilde{a}_{k-2}^s},\widetilde{z}_{k-1}^{s})}
\widetilde{q}^{\Delta}(\widetilde{z}_k^i|\widetilde{z}_{k-1}^{\widetilde{a}_{k-1}^i})
\right\}.
$$
Therefore,  applying \eqref{eq:prf2} and (A\ref{ass:1}) for any fixed $\widetilde{z}_{1:n}^{1:N},\widetilde{a}_{1:n-1}^{1:N}$ we have that
$$
\lim_{\Delta\downarrow 0} \mathbb{P}^{\Delta}_n(\widetilde{z}_{1:n}^{1:N},\widetilde{a}_{1:n-1}^{1:N}) = 
$$
$$
\left\{\prod_{i=1}^N p_1(\overline{z}_1^i)q_1^{\star}(z_1^i|x_0)\right\}\left\{
\prod_{k=2}^n \prod_{i=1}^N \frac{w_{k-1}^{\star}(z_{k-2}^{\widetilde{a}_{k-2}^i},z_{k-1}^{\widetilde{a}_{k-1}^i})}{
\sum_{s=1}^N
w_{k-1}^{\star}(z_{k-2}^{\widetilde{a}_{k-2}^s},z_{k-1}^{s})}
p_n(\overline{z}_k^i) q^{\star}(z_k^i|z_{k-1}^{\widetilde{a}_{k-1}^i})
\right\}.
$$
So by the dominated convergence theorem we have that
$$
\lim_{\Delta \downarrow 0}\mathbb{E}^{\Delta}\left[\left|\pi_n^{\Delta,N}(\varphi)-\pi_n^{\Delta}(\varphi)\right|^r\right] = 
$$
$$
\sum_{\widetilde{\mathbf{a}}_{1:n-1}^{1:N}\in\{1,\dots,N\}^{n-1}} \int_{\mathbb{R}^{Nnd_x}}
\left|
\frac{\sum_{i=1}^N \varphi(z_n^i)
w_{n}^{\star}(z_{n-1}^{\widetilde{a}_{n-1}^i},z_{n}^{i})}
{\sum_{i=1}^N
w_{n}^{\star}(z_{n-1}^{\widetilde{a}_{n-1}^i},z_{n}^{i})}
-\pi_n^{\star}(\varphi)
\right|^{r} \times
$$
$$
\left\{\prod_{i=1}^N p_1(\overline{z}_1^i)q_1^{\star}(z_1^i|x_0)\right\}\left\{
\prod_{k=2}^n \prod_{i=1}^N \frac{w_{k-1}^{\star}(z_{k-2}^{\widetilde{a}_{k-2}^i},z_{k-1}^{\widetilde{a}_{k-1}^i})}{
\sum_{s=1}^N
w_{k-1}^{\star}(z_{k-2}^{\widetilde{a}_{k-2}^s},z_{k-1}^{s})}
p_n(\overline{z}_k^i) q^{\star}(z_k^i|z_{k-1}^{\widetilde{a}_{k-1}^i})
\right\}
\bigotimes_{k=1}^n\bigotimes_{i=1}^N \lambda_{d_x}(d\widetilde{z}^i_k) = 
$$
$$
\sum_{\widetilde{\mathbf{a}}_{1:n-1}^{1:N}\in\{1,\dots,N\}^{n-1}} \int_{\mathbb{R}^{Nn(d_x-d_y)}}
\left|
\frac{\sum_{i=1}^N \varphi(z_n^i)
w_{n}^{\star}(z_{n-1}^{\widetilde{a}_{n-1}^i},z_{n}^{i})}
{\sum_{i=1}^N
w_{n}^{\star}(z_{n-1}^{\widetilde{a}_{n-1}^i},z_{n}^{i})}
-\pi_n^{\star}(\varphi)
\right|^{r} \times
$$
$$
\left\{\prod_{i=1}^N q_1^{\star}(z_1^i|x_0)\right\}\left\{
\prod_{k=2}^n \prod_{i=1}^N \frac{w_{k-1}^{\star}(z_{k-2}^{\widetilde{a}_{k-2}^i},z_{k-1}^{\widetilde{a}_{k-1}^i})}{
\sum_{s=1}^N
w_{k-1}^{\star}(z_{k-2}^{\widetilde{a}_{k-2}^s},z_{k-1}^{s})}
 q^{\star}(z_k^i|z_{k-1}^{\widetilde{a}_{k-1}^i})
\right\}
\bigotimes_{k=1}^n\bigotimes_{i=1}^N \lambda_{d_x}(dz^i_k) 
$$
where $\widetilde{\mathbf{a}}_{1:n-1}^{1:N} = (\widetilde{a}_1^1,\dots,\widetilde{a}_1^N,\dots,\widetilde{a}_{n-1}^1,\dots,\widetilde{a}_{n-1}^N)$.  We recognize the last expression as the expectation w.r.t.~the law of the degenerate noise PF; i.e.~we have proved that
$$
\lim_{\Delta \downarrow 0}\mathbb{E}^{\Delta}\left[\left|\pi_n^{\Delta,N}(\varphi)-\pi_n^{\Delta}(\varphi)\right|^r\right]^{1/r} = 
\mathbb{E}^{\star}\left[\left|\pi_n^{\star,N}(\varphi)-\pi_n^{\star}(\varphi)\right|^r\right]^{1/r}.
$$
\end{proof}

\subsection{Proof of Proposition \ref{prop:prop}}\label{app:inc}

\begin{proof}
The results under the natural parameterization are simple consequences of the choice of the proposal.  For the low-noise case using simple linear algebra one can verify that
$$
p_n(u_n^{\Delta}(\widetilde{z}_n,\epsilon))
f_n(u_n^{\Delta}(\widetilde{z}_n,x)|u_{n-1}^{\Delta}(\widetilde{z}_{n-1},x)) \propto
\exp\left\{
-\frac{1}{2}z_n^{\top}\widetilde{\Omega}_n^{-1}z_n + \widetilde{\mu}_n^{\top}z_n
\right\}.
$$  
From there one has the optimal proposal and the normalizing constant,  which is (up-to a constant) $w_n^{\Delta,\text{opt}}(\widetilde{z}_{n-1})$ can be obtained by completing the square and using Gaussian integrals.  For
the degenerate case we have similarly that
$$
f_n(u_n^{\star}(z_n)|u_{n-1}^{\star}(z_{n-1})) \propto
\exp\left\{
-\frac{1}{2}z_n^{\top}(\Omega_n^{\star})^{-1}z_n + (\mu_n^{\star})^{\top}z_n
\right\}
$$  
and the proof is similar to the low-noise case. The final limit is a direct consequence of \eqref{eq:prf1}
and the dominated convergence theorem.
\end{proof}


\subsection{Verification of Assumption (A\ref{ass:1}) in the Gaussian case}\label{app:a1_gaussian}

\begin{proof}

We consider the model
\[
Y_n = A_n X_n + \Delta^{1/2}\varepsilon_n,
\qquad
X_n = f_n(X_{n-1}) + \nu_n,
\]
where $\varepsilon_n \sim \mathcal N(0,\Sigma_n)$ and $\nu_n \sim \mathcal N(0,\Omega_n)$ are independent.
Let $y_n$ be fixed.

Let $x_n^\star \in \mathbb R^{d_x}$ satisfy $A_n x_n^\star = y_n$, and let
$V_n \in \mathbb R^{d_x \times (d_x-d_y)}$ have orthonormal columns spanning $\ker(A_n)$.
Define the coordinate map
\[
(x_n,\varepsilon_n)
=
\begin{pmatrix}
x_n^\star \\ 0
\end{pmatrix}
+
V_n^\Delta
\tilde z_n,
\qquad
\tilde z_n =
\begin{pmatrix}
z_n \\ \bar z_n
\end{pmatrix}
\in \mathbb R^{d_x},
\]
where
\[
V_n^\Delta =
\begin{pmatrix}
V_n & \Delta^{1/2} A_n^\top \\
0 & - A A^\top%-I_{d_y}
\end{pmatrix}.
\]
This implies
\[
x_n = x_n^\star + V_n z_n + \Delta^{1/2}A_n^\top \bar z_n,
\qquad
\varepsilon_n = - A_n A_n^\top \bar z_n.
\]

The optimal proposal density in the noise parametrisation is
\[
\tilde q_n^\Delta(\tilde z_n \mid \tilde z_{n-1})
\;\propto\;
p(\varepsilon_n)\, f(x_n \mid x_{n-1}).
\]
Substituting the parametrisation and writing
$m_n := f_n(x_{n-1}) - x_n^\star$, we obtain
\[
\tilde q_n^\Delta(\tilde z_n \mid \tilde z_{n-1})
=
\mathcal N_{d_x}(\tilde z_n;\mu_\Delta,\Sigma_\Delta),
\]
where the precision matrix $\Lambda_\Delta = \Sigma_\Delta^{-1}$ is
\[
\Lambda_\Delta =
\begin{pmatrix}
V_n^\top\Omega_n^{-1}V_n
&
\Delta^{1/2}V_n^\top\Omega_n^{-1}A_n^\top
\\[1mm]
\Delta^{1/2}A_n\Omega_n^{-1}V_n
&
\Sigma_n^{-1} + \Delta A_n\Omega_n^{-1}A_n^\top
\end{pmatrix},
\]
and the linear term equals
\[
\begin{pmatrix}
V_n^\top\Omega_n^{-1}m_n
\\[1mm]
\Delta^{1/2}A_n\Omega_n^{-1}m_n
\end{pmatrix}.
\]
Since $V_n^\top\Omega_n^{-1}V_n$ and $\Sigma_n^{-1}$ are positive definite,
block matrix inversion yields
\[
\Sigma_\Delta
\;\xrightarrow{\Delta \to 0}\;
\begin{pmatrix}
(V_n^\top\Omega_n^{-1}V_n)^{-1} & 0
\\[1mm]
0 & \Sigma_n
\end{pmatrix}.
\]
Multiplying the limiting covariance with the linear term gives
\[
\mu_\Delta
\;\xrightarrow{\Delta \to 0}\;
\begin{pmatrix}
(V_n^\top\Omega_n^{-1}V_n)^{-1}V_n^\top\Omega_n^{-1}m_n
\\[1mm]
0
\end{pmatrix}
=: \begin{pmatrix} \mu_n^\star \\ 0 \end{pmatrix}.
\]

Hence, as $\Delta \downarrow 0$,
\[
\tilde q_n^\Delta(z_n,\bar z_n \mid \tilde z_{n-1})
\;\Rightarrow\;
\mathcal N(z_n;\mu_n^\star,(V_n^\top\Omega_n^{-1}V_n)^{-1})\,
\mathcal N(\bar z_n;0,\Sigma_n).
\]
Writing $q_n^\star(z_n\mid z_{n-1})$ for the first factor and
$p_n(\bar z_n)$ for the second, we obtain
\[
\lim_{\Delta\downarrow 0}
\tilde q_n^\Delta(\tilde z_n \mid \tilde z_{n-1})
=
p_n(\bar z_n)\, q_n^\star(z_n\mid z_{n-1}),
\]
which verifies Assumption (A\ref{ass:1}).
\end{proof}


\subsection{Verification of Assumption (A\ref{ass:1}) for conditionally optimal $t$--noise proposal}\label{app:a1_tdist}
\begin{proof}

We consider the nonlinear state--space model used in the numerical example:
\[
Y_n = A_n X_n + \Delta^{1/2}\varepsilon_n,
\qquad
X_n = f(X_{n-1}) + \nu_n,
\]
where $\varepsilon_n \sim \mathcal N(0,\Sigma_n)$ and the transition noise
$\nu_n$ follows a multivariate $t$ distribution
\[
\nu_n \sim t_{d_x}(0,\Omega_n,\vartheta).
\]
The $t$--distributed transition admits the Gaussian scale mixture representation
\[
\nu_n \mid \omega_n \sim \mathcal N\!\left(0,\frac{\Omega_n}{\omega_n}\right),
\qquad
\omega_n \sim \mathrm{Gamma}\!\left(\frac{\vartheta}{2},\frac{\vartheta}{2}\right),
\]
with $\omega_n$ independent across $n$ and independent of all other variables.
Conditionally on $\omega_n$, the state equation becomes
\[
X_n = f(X_{n-1}) + \tilde\nu_n,
\qquad
\tilde\nu_n \sim \mathcal N\!\left(0,\frac{\Omega_n}{\omega_n}\right).
\]

As in Section \eqref{sec:noise_add}, introduce coordinates $\tilde z_n=(z_n,\bar z_n)\in\mathbb R^{d_x}$
via
\[
(x_n,\varepsilon_n)
=
\begin{pmatrix}
x_n^\star \\ 0
\end{pmatrix}
+
V_n^\Delta \tilde z_n,
\]
where $A_n x_n^\star=y_n$ and $V_n^\Delta$ is the lifted basis defined in the paper.
This parametrisation is independent of $\omega_n$.
Conditionally on $\omega_n$, the proposal used in the numerics is the optimal
Gaussian proposal in the noise parametrisation, namely
\[
\tilde q_n^\Delta(\tilde z_n \mid \tilde z_{n-1},\omega_n)
\;\propto\;
p(\varepsilon_n)\,
f(x_n\mid x_{n-1},\omega_n),
\]
where
\[
f(x_n\mid x_{n-1},\omega_n)
=
\mathcal N\!\left(
x_n; f(x_{n-1}),\frac{\Omega_n}{\omega_n}
\right).
\]
For each fixed $\omega_n>0$, this is exactly the Gaussian additive--noise setting
of \eqref{app:a1_gaussian}, with covariance matrix $\Omega_n/\omega_n$ in place of $\Omega_n$.
Hence, by the Gaussian case already established, we have for each fixed $\omega_n$
\[
\lim_{\Delta\downarrow 0}
\tilde q_n^\Delta(\tilde z_n \mid \tilde z_{n-1},\omega_n)
=
p_n(\bar z_n)\,
q_{n,\omega_n}^\star(z_n\mid z_{n-1}),
\]
where $q_{n,\omega_n}^\star$ is the optimal degenerate--noise proposal on the manifold
associated with covariance $\Omega_n/\omega_n$.

The actual proposal used in the algorithm samples $\omega_n$ from its prior
$\mathrm{Gamma}(\vartheta/2,\vartheta/2)$ and then samples
$\tilde z_n$ from the conditional proposal. Thus the marginal proposal density is
\[
\tilde q_n^\Delta(\tilde z_n \mid \tilde z_{n-1})
=
\int
\tilde q_n^\Delta(\tilde z_n \mid \tilde z_{n-1},\omega_n)\,
\pi(\omega_n)\, d\omega_n.
\]
By dominated convergence (the integrand is bounded by an integrable Gaussian
density uniformly in $\Delta$), we may pass the limit inside the integral to obtain
\[
\lim_{\Delta\downarrow 0}
\tilde q_n^\Delta(\tilde z_n \mid \tilde z_{n-1})
=
p_n(\bar z_n)
\int q_{n,\omega_n}^\star(z_n\mid z_{n-1})\,\pi(\omega_n)\,d\omega_n.
\]
Define
\[
q_n^\star(z_n\mid z_{n-1})
:=
\int q_{n,\omega_n}^\star(z_n\mid z_{n-1})\,\pi(\omega_n)\,d\omega_n,
\]
which is precisely the degenerate--noise proposal induced by the $t$--distributed
transition.

We have shown that
\[
\lim_{\Delta\downarrow 0}
\tilde q_n^\Delta(\tilde z_n \mid \tilde z_{n-1})
=
p_n(\bar z_n)\, q_n^\star(z_n\mid z_{n-1}),
\]
and therefore the conditionally optimal proposal used in the nonlinear
$t$--noise example satisfies Assumption (A\ref{ass:1}).
\end{proof}


\begin{thebibliography}{99}

\bibitem{diff_bridge}
{\sc Alvarez}, M.  \& {\sc Jasra}, A.~(2025). 
Unbiased Parameter Estimation of Partially Observed Diffusions using Diffusion Bridges.  arXiv preprint.


\bibitem{bain}
{\sc Bain}, A., \& {\sc Crisan}, D.~(2009). \emph{Fundamentals of Stochastic Filtering}. Springer: New York.

\bibitem{app1}
{\sc Barczyk}, M. {\sc Bonnabel,} S., {\sc Deschaud}, J. \& {\sc Goulette},  F.~(2015).
Invariant EKF design for scan matching-aided localization. \emph{IEEE Trans. Control Syst. Technol.}, 
{\bf 23}, 2440-2448.


\bibitem{beskos}
{\sc Beskos}, A., {\sc Crisan}, D., {\sc Jasra}, A., {\sc Kantas}, N. \& {\sc Ruzayqat}, H.~(2021).
Score-Based parameter estimation for a class of continuous-time state space models.
\emph{SIAM J. Sci. Comp.}, {\bf 43},  A2555-A2580.


\bibitem{bhar}
{\sc Bharath}, K., {\sc Lewis}, A., {\sc Sharma},  A. {\sc Tretyakov},  M.~(2025).  Sampling and Estimation on Manifolds using the Langevin Diffusion. \emph{J. Mach. Learn. Res.} (to appear). 


\bibitem{brigo}
{\sc Brigo}, D.~(1996). New results on the Gaussian projection filter with small observation noise. \emph{Systems \& control letters}, {\bf 28} (5), 273-279.



\bibitem{bierkens}
{\sc Bierkens}, J.,  \textsc{van der Meulen}, F., \& \textsc{Schauer}, M.~(2020).  Simulation of elliptic and hypoelliptic conditional diffusions.  \emph{Adv.  Appl.  Probab.},  {\bf 52},  173-212.

\bibitem{byrne}
{\sc Byrne}, S. \& {\sc Girolami},  M.~(2013). Geodesic Monte Carlo on embedded manifolds.  \emph{Scand.  J.  Stat.}
{\bf 40},  825--845.

\bibitem{cappe}
{\sc Capp\'e}, O., {\sc Ryden}, T, \& {\sc Moulines}, \'E.~(2005). \emph{Inference
in Hidden Markov Models}. Springer: New York.

\bibitem{delmoral}
{\sc Del Moral}, P.~(2004).
\emph{Feynman-Kac Formulae: Genealogical and Interacting Particle Systems with Applications}. Springer: New York.



\bibitem{graham_st}
{\sc Graham,} M.,  \& {\sc Storkey},  A.~(2017). Asymptotically exact inference in differentiable generative models.
\emph{Elec. J. Stat.}, {\bf 11},  5105--5164.

\bibitem{graham}
{\sc Graham,} M., {\sc Thiery,} A. \& {\sc Beskos}, A.~(2022).
Manifold Markov chain Monte Carlo methods for Bayesian inference in diffusion models,  \emph{J.~R.~Statist.~Soc.~B}, {\bf 84},  1229-1256.


\bibitem{hijab}
{\sc Hijab}, O.~(1981) Asymptotic Bayesian Estimation of a First Order Equation with Small Diffusion. \emph{Ann. Probab.}, {\bf 12} (3) 890 -- 902.  

\bibitem{app2}
{\sc Hua},  A.,  {\sc Ducard}, G.,  {\sc Hamel}, T.  {\sc Mahony},  R. \& {\sc Rudin}, K.~(2014).
Implementation of a nonlinear attitude estimator for aerial robotic vehicles.
\emph{IEEE Trans. Control Syst. Technol.},  {\bf 22},  201-213.


\bibitem{non_synch}
{\sc Jasra}, A.,  {\sc Kamatani}, K. \& {\sc Wu},  A.~(2026). Bayesian Inference for non-synchronously observed diffusions.  \emph{SIAM/ASA JUQ},  (to appear).

\bibitem{pf_ns}
{\sc Kantas}, N., {\sc Beskos}, A.  \& {\sc Jasra}, A.~(2014). SMC for inverse problems: a case study for Navier Stokes. \emph{SIAM/ASA JUQ},  {\bf 2}, 464-489.



\bibitem{katzur}
{\sc Katzur}, R., {\sc Bobrovsky}, B.Z., and {\sc Schuss}, Z. (1984) Asymptotic analysis
of the optimal filtering problem for one-dimensional diffusions
measured in a low noise channel, Part I and II, \emph{SIAM J. Appl.
Math.}, {\bf 44}, pp. 591-604, and pp. 1176-1191.

\bibitem{kutoyants}
{\sc Kutoyants}, Y. A.~(2025).  Extended adaptive Kalman filter with low noise observations. \emph{Probability, Uncertainty and Quantitative Risk.}, {\bf 10}(4), 443--470. 



\bibitem{milheiro}
{\sc Milheiro De Oliveira}, P.~(1994) Approximate filters for a nonlinear discrete time filtering problem with small observation noise, \emph{Stochastics and Stochastic Reports},  {\bf 46}(1-2): 1--24.


\bibitem{pardoux}
{\sc Pardoux}, E., \& {\sc Zeitouni}, O.~(2004). Quenched large deviations for one dimensional nonlinear filtering. \emph{SIAM Journal  on Control and Optimization}, {\bf 43} (4), 1272--1297.


\bibitem{picard1}
{\sc Picard}, J.~(1986)  Nonlinear filtering of one-dimensional diffusion in the case of a high signal-to-noise ratio, \emph{SIAM Journal on Applied Mathematics}, {\bf 46}(6): 1098--1125.

\bibitem{picard2}
{\sc Picard}, J.~(1991)  Efficiency of the extended Kalman filter for nonlinear systems with small noise, \emph{SIAM Journal on Applied Mathematics}, {\bf 51}(4): 843--885.

\bibitem{pf_sns}
{\sc Pons Llopis}, F., {\sc Kantas}, N., {\sc Beskos}, A.  \& {\sc Jasra}, A.~(2018). Particle filtering for stochastic Navier Stokes with additive noise.  \emph{SIAM J. Sci. Comp. }, {\bf 40},  A1544-A1565.



\bibitem{schauer}
{\sc Schauer}, M., {\sc van der Meulen}, F. \& {\sc van Zanten}, H.~(2017). Guided proposals for simulating multi-dimensional diffusion bridges. \emph{Bernoulli}, {\bf 23}, 2917--2950.

\bibitem{stanton}
{\sc Stanton},   C.  \& {\sc Beskos},  A.~(2024).  Particle based inference for continuous-discrete state space models.  arXiv preprint.

\bibitem{vd_meulen_guided_mcmc}
\textsc{van der Meulen}, F., \& \textsc{Schauer}, M.~(2017). Bayesian estimation of discretely observed multi-dimensional diffusion processes using guided proposals. \emph{Elec. J. Stat.}, {\bf 11}, 2358-2396. 

\bibitem{shouto}
\textsc{Yonekura}, S., \& \textsc{Beskos}, A.~(2022). Online Smoothing for Diffusion Processes Observed with Noise. \emph{J. Comp. Graph. Stat.}, {\bf 31}(4), 1344--1360. 

\bibitem{zappa}
{\sc Zappa},  E.,  {\sc Holmes-Cerfon},  A. \& {\sc Goodman},  J.~(2018).  Monte Carlo on manifolds: sampling densities and integrating functions.  \emph{Comm.  Pure Appl.  Math.}, {\bf 71},  2609-2647.

\bibitem{zhang}
{\sc Zhang},  C.,  {\sc Taghvaei},  A., {\sc Mehta}, P.~(2018).  Feedback Particle Filter on Riemannian Manifolds and Matrix Lie Groups.  \emph{IEEE Trans.  Aut. Contr.},  {\bf 63},  2465--2480.

\bibitem{zhu}
{\sc Zhumekenov},  A.,  {\sc Beskos}, A., {\sc Crisan}, D.,  {\sc Jasra}, A.,  \& {\sc Kantas}, N.~(2025).
Sequential Markov chain Monte Carlo for Filtering of State-Space Models with Low or Degenerate Observation Noise.  arXiv preprint.

\end{thebibliography}



\end{document}

